%%%%%%%%%%%%%%%%%%%%%%%%%%%%%%%%%%%%%%%%%%%%%%%%%%%%%%%%%%%%%%%%%%%%%%
% Lớp 12 | Các số đặc trưng ... ghép nhóm | Khoảng biến thiên, khoảng tứ phân vị của mẫu số liệu ghép nhóm
% Dạng 2D3?1-3: Tìm khoảng tứ phân vị 
% Số câu: 45  (sinh tự động bởi tach_cau.py; từ nay thầy sửa trực tiếp file này)
%%%%%%%%%%%%%%%%%%%%%%%%%%%%%%%%%%%%%%%%%%%%%%%%%%%%%%%%%%%%%%%%%%%%%%

%%%%%%%%%%%%%%%% TRẮC NGHIỆM 4 PHƯƠNG ÁN %%%%%%%%%%%%%%%%

% nguồn: Lop_12_Dai_so_Chuong_3
\begin{ex}%[2D3N1-3]%[ma:67b65b2c]
	Một  mẫu số liệu ghép nhóm có tứ phân vị thứ nhất, thứ hai, thứ ba lần lượt là $Q_1$, $Q_2$, $Q_3$. Khoảng tứ phân vị của mẫu số liệu đó bằng
	 \choice
		{$Q_2-Q_1$} 
		 {\True $Q_3-Q_1$}
		 {$Q_3-Q_2$}
		{$Q_3+Q_1-Q_2$} 
	\loigiai{Khoảng tứ phân vị của mẫu số liệu đó  là $Q_3-Q_1$.}
\end{ex}

% nguồn: Lop_12_Dai_so_Chuong_3
\begin{ex}%[2D3H1-3]%[ma:96427650]
	Trung tâm ngoại ngữ thống kê bảng điểm môn Tiếng Anh của một khóa học trong bảng bên  dưới
	\begin{center}
		\begin{tabular}{|c|c|c|c|c|c|}
			\hline 
			Điểm&$[0;2)$&$[2;4)$&$[4;6)$&$[6;8)$&$[8;10)$\\ 
			 \hline
			Học  viên&$10$&$30$&$55$&$42$&$9$\\
			\hline 
		\end{tabular}
	\end{center}
	Khoảng  tứ phân vị của mẫu số liệu ghép nhóm này là
	\choice 
		{\True$2{,}92$} 
		{$2{,}93$} 
		 {$3{,}92$}
		{$3{,}93$} 
	 \loigiai{
		Ta  có $n=146$.\\
		Tứ phân vị thứ nhất của mẫu số liệu gốc là $x_{37}\in[2;4)$. Tứ phân vị thứ ba của mẫu liệu  gốc là $x_{110}\in[6;8)$.\\
		Tứ phân vị thứ  nhất của mẫu số liệu ghép nhóm trên là
		 $$Q_1=2+\dfrac{\tfrac{1\cdot146}{4}-10}{30}\cdot(4-2)=\dfrac{113}{30}.$$
		Tứ phân vị thứ ba của  mẫu số liệu ghép nhóm trên là
		 $$Q_3=6+\dfrac{\tfrac{3\cdot146}{4}-(10+30+55)}{42}\cdot(8-6)=\dfrac{281}{42}.$$
		Khoảng tứ phân vị  $\Delta Q=Q_3-Q_1=\dfrac{281}{42}-\dfrac{113}{30}\approx2{,}92$.
	} 
\end{ex}

% nguồn: Lop_12_Dai_so_Chuong_3
\begin{ex}%[2D3H1-3]%[ma:223d04d9]
	Mỗi ngày bác Hương  đều đi bộ để rèn luyện sức khoẻ. Quãng đường đi bộ mỗi ngày (đơn vị: km) của bác Hương trong $20$ ngày được thống kê lại ở bảng sau
	\begin{center}
		\begin{tabular}{|c|c|c|c|c|c|}
			\hline \begin{tabular}{c}
			 Quãng đường \\
			 km
		\end{tabular} & {$[2{,}7 \,; \,3{,}0)$} & {$[3{,}0 \,; \,3{,}3)$} & {$[3{,}3 \,;\, 3{,}6)$} & {$[3{,}6 \,; \,3{,}9)$} & {$[3{,}9 \,; \,4{,}2)$} \\
		\hline Số ngày & $3$ & $6$ & $5$ & $4$ & $2$  \\
		\hline 
	\end{tabular}
\end{center}
Khoảng tứ phân vị của mẫu  số liệu ghép nhóm là
\choice 
	 {$0{,}9$}
	 {$0{,}975$}
	{$0{,}5$} 
	{\True  $0{,}575$}
 \loigiai{
	Cỡ mẫu $n=20$. \\ 
	Gọi $x_{1} $; $x_{2}$; $\ldots $; $x_{20}$ là mẫu số liệu gốc về quãng đường đi  bộ mỗi ngày của bác Hương trong $20$ ngày được xếp theo thứ tự không giảm.\\
	Ta có $x_{1} , \ldots , x_{3} \in[2{,}7; 3{,}0) $; $x_{4} , \ldots , x_{9} \in[3{,}0 \,;\, 3{,}3) $; $x_{10}  ,\ldots , x_{14} \in[3{,}3 \,; \,3{,}6)$; $x_{15} , \ldots , x_{18} \in[3{,}6 \,; \,3{,}9) $; $x_{19} , x_{20} \in[3{,}9 \,; \,4{,}2)$.\\
	Tứ phân vị thứ nhất của mẫu số liệu gốc  là $\dfrac{1}{2}(x_{5}+x_{6}) \in[3{,}0 \,; \,3{,}3)$.\\
	Do đó, tứ phân vị thứ nhất của mẫu số liệu ghép nhóm  là $Q_{1}=3{,}0+\dfrac{\dfrac{20}{4}-3}{6}\cdot(3{,}3-3{,}0)=3{,}1$.\\
	Tứ phân vị thứ ba của mẫu số liệu gốc là  $\dfrac{1}{2}(x_{15}+x_{16}) \in[3{,}6 \,; \,3{,}9)$.\\
	Do  đó, tứ phân vị thứ ba của mẫu số liệu ghép nhóm là
	$$ 
	 Q_{3}=3{,}6+\dfrac{\dfrac{3\cdot 20}{4}-(3+6+5)}{4}\cdot(3{,}9-3{,}6)=3{,}675.
	 $$
Khoảng tứ phân  vị của mẫu số liệu ghép nhóm là $\Delta_{Q}=Q_{3}-Q_{1}=0{,}575$.}
\end{ex}

% nguồn: Lop_12_Dai_so_Chuong_3
\begin{ex}%[2D3H1-3]%[ma:ebb70be6]
	Tại một gara ô tô, người ta ghi lại thời gian sử dụng của một số chiếc ô tô đến sữa chữa tại đây (đơn  vị: năm). Kết quả được cho trong bảng sau đây.
	\begin{center}
		\begin{tabular}{|c|c|c|c|c|c|c|}
			\hline Nhóm & $[4; 6)$ & $[6; 8)$ & $[8; 10)$  & $[10; 12)$ & $[12; 14)$ & $[14; 16)$ \\
			\hline Tần  số & $10$ & $14$ & $19$ & $23$ & $27$ & $32$ \\
			 \hline
		\end{tabular}
	\end{center}
	Khoảng tứ phân vị của mẫu số liệu ghép  nhóm trên (làm tròn các kết quả phép tính đến hàng phần trăm) là
	 \choice
		{\True  $5{,}29$}
		{$4{,}26$} 
		 {$6$}
		{$5{,}07$} 
	\loigiai{ 
		Cỡ  mẫu $n=125$.\\
		 Tứ phân vị thứ nhất của mẫu số liệu là $Q_1=\dfrac{x_{31}+x_{32}}{2}\in [8;10)$. Lúc đó\[Q_1=8+\dfrac{\dfrac{1\cdot 125}{4}-24}{19}\cdot (10-8)\approx 8{,}76.\]
		Tứ phân vị thứ ba của mẫu số liệu là $Q_3=\dfrac{x_{95}+x_{96}}{2}\in  [14;16)$. Lúc đó
		\[Q_3=14+\dfrac{\dfrac{3\cdot  125}{4}-93}{32}\cdot (16-14)\approx 14{,}05.\]
		Vậy khoảng tứ phân vị của mẫu số  liệu ghép nhóm trên (làm tròn đến hàng phần trăm) là $\Delta_Q=Q_3-Q_1\approx 5{,}29$.
	} 
\end{ex}

% nguồn: Lop_12_Dai_so_Chuong_3
\begin{ex}%[2D3H1-3]%[ma:a6f5e787]
	Biểu đồ  dưới đây biểu diễn số lượng khách hàng đặt bàn qua hình thức trực tuyến mỗi ngày trong quý $III$ năm $2022$ của một nhà hàng. Cột thứ nhất biểu diễn số ngày có từ $1$ đến $6$ lượt đặt bàn, cột thứ hai biễu diễn số ngày có từ $6$ đến $11$ lượt đặt bàn;...
	\begin{center}
		\begin{tikzpicture}[>=stealth,declare function={ox=-0.5; kc=1; rc=1.0; kcy=0.9;}]
				 \foreach[count=\i] \yt in {5,10,15,20,25,30,35}{\draw[dashed] (0,{\i*kcy}) node[left,text=black]{$\yt$}--+(7,0);}
				\foreach[count=\i] \yt/\xt in {2.8/1,6/6,5/11,3.8/16,1/21,0/26}{\fill[cyan] ({ox+\i*kc},0) rectangle +(rc,{\yt*kcy}); \path ({ox+\i*kc+0.5*rc-0.2},0) node[below left]{\xt}; \path ({ox+\i*kc+0.5*rc},{\yt*kcy});} 
				\foreach[count=\i] \yt/\xt in {2.8/I,6/6,5/11,3.8/16,1/21,0/26}{\draw ({ox+\i*kc},0) rectangle +(rc,{\yt*kcy}); \path  ({ox+\i*kc+0.5*rc},{\yt*kcy});}
				\foreach[count=\i] \yt/\gt in {2.8/14,6/30,5/25,3.8/18,1/5}{\draw ({ox+\i*kc},0) rectangle +(rc,{\yt*kcy});  \path ({ox+\i*kc+0.5*rc},{\yt*kcy}) node[above]{$\gt$};}
				 \draw[->] (0,0)--(7.5,0) node[below]{Số lượt đặt bàn};
				 \draw[->] (0,0)--(0,7) node[left,anchor=east]{Số ngày};
			\end{tikzpicture}
	\end{center}
	 Khoảng tứ phân vị của mẫu số liệu ghép nhóm cho bởi biểu đồ trên
	\choice 
		 {$9{,}5$}
		{\True $8{,}5$} 
		{$10{,}5$} 
		 {$7{,}5$}
	\loigiai{ 
		Dựa vào  biểu đồ ta lập được bảng ghép nhóm như sau
		\begin{center}
			\begin{tabular}{|c|c|c|c|c|c|}
				\hline Lượt  đặt bàn & {$[1; 6)$} & {$[6; 11)$} & {$[11; 16)$} & {$[16; 21)$} & {$[21; 26)$} \\
				\hline Số ngày & $14$ & $30$ & $25$  & $18$ & $5$ \\
				 \hline
			\end{tabular}
		\end{center}
		Cỡ mẫu $n=92$ và gọi  $x_{1}$, $x_{2}$, $\ldots$, $x_{92}$ là mẫu số liệu đã cho.\\
		Ta có $x_{1}, \ldots, x_{14} \in[1;  6) $; $x_{45}, \ldots, x_{69} \in[11; 16)$; $x_{70}, \ldots , x_{87} \in[16; 21) $; $ x_{88}, \ldots, x_{92} \in[21; 26)$.\\
		Tứ phân vị  thứ nhất của mẫu số liệu là $\dfrac{x_{23}+x_{24}}{2} \in[6; 11)$. \\
		Do đó tứ phân vị thứ nhất của mẫu số liệu  là \[Q_{1}=6+\dfrac{\dfrac{92}{4}-14}{30}(11-6)=7{,}5.\]
		Tứ phân vị thứ ba của mẫu số liệu là $\dfrac{x_{69}+x_{70}}{2} \in[11; 16)$  và $x_{70} \in[16; 21)$.\\
		Do đó  tứ phân bị thứ ba của mẫu số liệu là $Q_{3}=16$.\\
	Vậy khoảng tứ phân vị của mẫu số liệu  là $\Delta_{Q}=Q_{3}-Q_{1}=16-7{,}5=8{,}5$.}
\end{ex}

% nguồn: Lop_12_Dai_so_Chuong_3
\begin{ex}%[2D3H1-3]%[ma:648b7111]
	Bảng sau đây biểu diễn  mẫu số liệu ghép nhóm về cân nặng của một số quả bưởi da xanh ở một khu vực (đơn vị: kg).
	\begin{center}
		\begin{tabular}{|c|c|c|c|c|c|}
			\hline Nhóm & $[1{,}2; 1{,}3)$ & $[1{,}3; 1{,}4)$ & $[1{,}4; 1{,}5)$ & $[1{,}5; 1{,}6)$ &  $[1{,}6; 1{,}7)$ \\
			\hline Tần  số & $14$ & $40$ & $13$ & $10$ & $3$ \\
			\hline 
		\end{tabular}
	\end{center}
	Khoảng tứ  phân vị của mẫu số liệu ghép nhóm trên (làm tròn đến hàng phần trăm) là
	 \choice
		{\True  $0{,}13$}
		 {$0{,}10$}
		 {$0{,}20$}
		{$0{,}14$} 
	 \loigiai{
		Tứ phân vị thứ nhất là $Q_1=\dfrac{x_{20}+x_{21}}{2}\in [1{,}3;1{,}4)$ và\[Q_1=1{,}3+\dfrac{\dfrac{1\cdot 180}{4}-14}{40}\cdot  (1{,}4-1{,}3)=1{,}315.\]
		Tứ phân vị thứ ba là $Q_3= 
		 \dfrac{x_{60}+x_{61}}{2}\in [1{,}4;1{,}5)$ và
		 \[Q_3=1{,}4+\dfrac{\dfrac{3\cdot 180}{4}-54}{13}\cdot (1{,}5-1{,}4)\approx 1{,}446.\]
		Vậy khoảng  tứ phân vị của mẫu số liệu trên là $\Delta_Q=Q_3-Q_1\approx 0{,}131$.
	 }
\end{ex}

% nguồn: Lop_12_Dai_so_Chuong_3
\begin{ex}%[2D3H1-3]%[ma:d4ffefb7]
	Điểm trung bình môn  Toán của một số học sinh khối $12$ được ghi lại ở bảng sau
	\begin{center}
		\begin{tabular}{|l|c|c|c|c|c|}
			\hline Nhóm  &{$[5;6)$}&{$[6;7)$}&{$[7;8)$}&{$[8;9)$}&{$[9;10)$}\\
			\hline Tần số & $10$ & $14$ & $19$ &  $32$ & $25$ \\
			\hline 
		\end{tabular}
	\end{center}
	Khoảng tứ phân vi của mẫu số liệu ghép nhóm trên  (làm tròn đến hàng phần trăm) là
	 \choice
		{$1{,}99$} 
		 {$1{,}98$}
		 {$1{,}96$}
		{\True $1{,}95$} 
	 \loigiai{Cỡ mẫu $n=100$.\\
		Gọi  $x_1; x_2; \ldots; x_{100}$ là mẫu số liệu gốc gồm điểm số của $100$ học sinh được xếp theo thứ tự không giảm.\\
		Ta có 
		\begin{listEX}[3]
			\item $x_1, \ldots,  x_{10} \in[5;6)$
			\item $x_{11}, \ldots,  x_{24} \in[6;7)$
			\item $x_{25}, \ldots,  x_{43} \in[7;8)$
			\item $x_{44}, \ldots, x_{75} \in[8;9)$ 
			\item $x_{76}, \ldots, x_{100}  \in[9;10)$.
		\end{listEX}
		Tứ phân vị thứ nhất của mẫu số liệu gốc là $\dfrac{1}{2}(x_{25}+x_{26})  \in[7;8)$.\\
		Do  đó, tứ phân vị thứ nhất của mẫu số liệu ghép nhóm là
		 \[Q_1=7+\dfrac{\dfrac{100}{4}-24}{19} \cdot(8-7)=\dfrac{134}{19}.\]
		Tứ phân vị thứ ba của mẫu số liệu gốc là  $\dfrac{1}{2}(x_{75}+x_{76})$.\\ Mà $x_{75} \in[8;9)$, $x_{76} \in[9;10)$.\\ Do đó, tứ phân vị thứ ba của mẫu số liệu ghép nhóm là
		\[Q_3=9.\] 
		Vậy khoảng tứ phân vị  của mẫu số liệu ghép nhóm là
		\[\Delta_Q=\dfrac{134}{19}-9  \approx 1{,}947.\]
	} 
\end{ex}

% nguồn: Lop_12_Dai_so_Chuong_3
\begin{ex}%[2D3H1-3]%[ma:31eff8d8]
	Thời gian (phút) truy cập Internet mỗi buổi tối  của một số học sinh được cho trong bảng sau
	\begin{center}
		\begin{tabular}{|c|c|c|c|c|c|}
			\hline \begin{tabular}{c}
			Thời  gian \\
			 (phút)
		\end{tabular} & {$[9{,}5; 12{,}5)$} & {$[12{,}5; 15{,}5)$} & {$[15{,}5; 18{,}5)$} & {$[18{,}5; 21{,}5)$} & {$[21{,}5; 24{,}5)$} \\
		\hline Số học sinh & $3$ & $12$ & $15$ & $24$ & $2$  \\
		\hline 
	\end{tabular}
\end{center}
Khoảng tứ phân vị của mẫu số liệu  ghép nhóm là
 \choice
	{$10{,}75$} 
	{\True  $4{,}75$}
	{$4{,}63$} 
	 {$4{,}38$}
\loigiai{ 
	Cỡ mẫu là $n=56$. Tứ phân vị thứ nhất $Q_{1}$ là  $\dfrac{x_{14}+x_{15}}{2}$. \\
	Do $x_{14},  x_{15}$ đều thuộc nhóm $[12{,}5; 15{,}5)$ nên nhóm này chứa $Q_{1}$. \\
	Do đó $p=2 $; $a_{2}=12{,}5 $; $m_{2}=12 $; $m_{1}=3$,  $a_{3}-a_{2}=3$ và ta có \[Q_{1}=12{,}5+\dfrac{\dfrac{56}{4}-3}{12} \cdot 3=15,25.\]
	Với tứ phân vị thứ  ba $Q_{3}$ là $\dfrac{x_{42}+x_{43}}{2}$. \\
	Do $x_{42}, x_{43}$ đều thuộc nhóm $[18{,}5; 21{,}5)$ nên  nhóm này chứa $Q_{3}$. \\
	 Do đó
	$p=4  $; $a_{4}=18{,}5 $; $m_{4}=24 $; $m_{1}+m_{2}+m_{3}=3+12+15=30 $; $ a_{5}-a_{4}=3$ và ta có
	 \[Q_{3}=18{,}5+\dfrac{\dfrac{3\cdot 56}{4}-30}{24} \cdot 3=20.\]
Khoảng tứ phân vị của mẫu số liệu ghép nhóm là $\Delta_{Q}=Q_{3}-Q_{1}=4{,}75$.} 
\end{ex}

% nguồn: Lop_12_Dai_so_Chuong_3
\begin{ex}%[2D3H1-3]%[ma:52bb7e27]
	Xét mẫu số liệu ghép nhóm có tứ phân vị thứ nhất, tứ phân vị thứ hai, tứ phân vị thứ ba lần lượt là $3$, $6$ và $8$. Khoảng tứ phân vị của mẫu số liệu ghép nhóm  đó bằng
	\choice 
		 {$3$}
		{$2$} 
		{\True $5$} 
		{$17$} 
	\loigiai{Khoảng tứ phân vị của mẫu số liệu ghép nhóm là hiệu của tứ phân vị thứ  ba và tứ phân vị thứ nhất. Do đó $\Delta_ Q=Q_3-Q_1=8-3=5$.
	} 
\end{ex}

% nguồn: Lop_12_Dai_so_Chuong_3
\begin{ex}%[2D3H1-3]%[ma:2708a01a]
	Bảng  sau ghi lại nhiệt độ không khí đo được lúc $12$ giờ trưa hằng ngày tại một trạm quan sát trong $1$ năm.
	\begin{center}
		\begin{tabular}{|l|c|c|c|c|c|}
			 \hline Nhóm &{$[12;16)$}&{$[16;20)$}&{$[20;24)$}&{$[24;28)$}&{$[28;32)$}\\
			\hline Tần số & $85$ & $63$ & $95$ & $87$ & $35$ \\ 
			 \hline
		\end{tabular}
	\end{center}
	Tứ phân vị thứ nhất của  mẫu số liệu ghép nhóm trên là
	\choice 
		{$18{,}397$} 
		{$12{,}397$} 
		{\True  $16{,}397$}
		{$16{,}099$} 
	\loigiai{Cỡ mẫu  $n=365$.\\
		Gọi $x_1; x_2; \ldots; x_{365}$ là mẫu số liệu gốc gồm nhiệt độ không khí đo được lúc 12 giờ trưa hằng ngày tại một trạm quan sát trong 1 năm được xếp  theo thứ tự không giảm.\\
		Ta có 
		\begin{listEX}[3]
			 \item $x_1, \ldots, x_{85} \in[12;16)$
			\item $x_{86}, \ldots, x_{148}  \in[16;20)$
			\item $x_{149},  \ldots, x_{243} \in[20;24)$
			\item  $x_{244}, \ldots, x_{330} \in[24;28)$
			\item  $x_{331} \ldots, x_{365}\in[28;32)$.
		\end{listEX}
		Tứ phân vị thứ nhất của mẫu số liệu gốc là $x_{91},x_{92}  \in[16;20)$.\\
		Do đó, tứ phân vị thứ nhất của mẫu số liệu ghép nhóm  là
		\[Q_1=16+\dfrac{\dfrac{365}{4}-85}{63} \cdot(20-16)\approx16{,}397.\] 
	 }
\end{ex}

% nguồn: Lop_12_Dai_so_Chuong_3
\begin{ex}%[2D3H1-3]%[ma:744145b9]
	Ghi lại tốc độ bóng trong $200$ lần giao bóng của một vận động viên  môn quần vợt cho kết quả như bảng bên.
	\begin{center}
		\begin{tabular}{|c|c|}
			\hline Tốc độ $v(\mathbf{k m} / \mathbf{h})$ & Số lần \\ 
			\hline  $150 \leq v<155$ & $18$ \\
			\hline $155 \leq v<160$ & $28$ \\ 
			\hline  $160 \leq v<165$ & $35$ \\
			\hline $165 \leq v<170$ &  $43$ \\
			\hline $170 \leq v<175$ & $41$  \\
			\hline  $175 \leq v<180$ & $35 $\\
			\hline 
		\end{tabular}
	\end{center}
	Khoảng tứ phân vị của  mẫu số liệu ghép nhóm là
	 \choice
		{\True $12{,}6$} 
		 {$11{,}5$}
		 {$14{,}3$}
		 {$16{,}8$}
	 \loigiai{
		 Cỡ mẫu là $n=200$.\\
		Tứ phân vị thứ nhất $Q_{1}$ là  $\dfrac{x_{50}+x_{51}}{2}$. \\
		Do $x_{50}, x_{51}$ đều thuộc nhóm $[160;  165)$ nên nhóm này chứa $Q_{1}$.
		Do đó $p=3 $;  $a_{3}=160 $; $m_{3}=35 $; $m_{1}+m_{2}=18+28=46$; $a_{4}-a_{3}=165-160=5$ và ta có
		\[Q_{1}=160+\dfrac{\dfrac{200}{4}-46}{35} \cdot  5 \approx 160{,}57.\]
		Tứ phân  vị thứ ba $Q_{3}$ là $\dfrac{x_{150}+x_{151}}{2}$. \\
		Do $x_{150}, x_{151}$ đều thuộc  nhóm $[170;175)$ nên nhóm này chứa $Q_{3}$. Do đó $p=5 $; $a_{5}=170$; $m_{5}=41 $; $m_{1}+m_{2}+m_{3}+m_{4}=18+28+35+43=124 $; $ a_{6}-a_{5}=175-170=5$ và ta có
		\[Q_{3}=170+\dfrac{\dfrac{3\cdot 200}{4}-124}{41} \cdot5 \approx 173{,}17.\] 
	Khoảng tứ phân vị của mẫu số liệu ghép nhóm là  $\Delta_{Q}=Q_{3}-Q_{1} \approx 12{,}6$.}
\end{ex}

% nguồn: Lop_12_Dai_so_Chuong_3
\begin{ex}%[2D3H1-3]%[ma:4dd7cdb1]
	Một vườn thú ghi lại tuổi thọ (đơn vị: năm) của $20$ con hổ và thu được kết quả như sau 
	\begin{center}
		\begin{tabular}{|c|c|c|c|c|c|}
			\hline Tuổi thọ & {$[14; 15)$} & {$[15; 16)$} & {$[16; 17)$} &  {$[17; 18)$} & {$[18; 19)$} \\
			\hline  Số con hổ & $1$ & $3$ & $8$ & $6$ & $2 $\\
			\hline 
		\end{tabular}
	\end{center}
	Nhóm  chứa tứ phân vị thứ nhất là
	\choice 
		{$[14;  15)$}
		 {$[15; 16)$}
		{\True  $(16;17)$}
		{$[17; 18)$} 
	 \loigiai{
		Ta có $\dfrac{n}{4}=\dfrac{20}{4}=5$ và $1+3<5<1+3+8$.\\ 
		Suy ra tứ phân vị  thứ nhất của mẫu số liệu thuộc nhóm $[16;17)$.
	} 
\end{ex}

% nguồn: Lop_12_Dai_so_Chuong_3
\begin{ex}%[2D3H1-3]%[ma:d279db0c]
	Bạn Chi rất thích nhảy hiện đại. Thời  gian tập nhảy mỗi ngày trong thời gian gần đây của bạn Chi được thống kê lại ở bảng sau
	\begin{center}
		\begin{tabular}{|c|c|c|c|c|c|}
			\hline \begin{tabular}{c}
			 Thời gian \\
			(phút) 
		\end{tabular} & {$[20; 25)$} & {$[25; 30)$} & {$[30; 35)$} & {$[35; 40)$} & {$[40; 45)$} \\
		\hline Số ngày &  $6$ & $6$ & $4$ & $1$ & $1$ \\
		 \hline
	\end{tabular}
\end{center}
Khoảng tứ phân vị của mẫu số liệu ghép nhóm  là
 \choice
	 {$23{,}75$}
	 {$27{,}5$}
	 {$31{,}88$}
	 {\True $8{,}125$}
\loigiai{ 
	Cỡ  mẫu $n=18$.\\
	Gọi $x_{1} $; $x_{2} $; $\ldots $; $x_{18}$ là mẫu số liệu gốc về thời gian tập nhảy mỗi ngày  của bạn Chi được xếp theo thứ tự không giảm. Ta có \\
	$x_{1} , \ldots , x_{6} \in[20; 25) $; $x_{7} , \ldots , x_{12} \in[25; 30) $; $x_{13}  , \ldots , x_{16} \in[30; 35) $; $x_{17} \in[35; 40) $; $x_{18} \in[40; 45)$.\\
	Tứ phân vị thứ nhất của mẫu số liệu gốc là $x_{5}  \in[20; 25)$. Do đó, tứ phân vị thứ nhất của mẫu số liệu ghép nhóm là \[Q_{1}=20+\dfrac{\dfrac{18}{4}}{6}\cdot (25-20)=23{,}75.\]
	Tứ phân vị thứ ba của mẫu số liệu gốc là $x_{14} \in[30; 35)$. Do đó, tứ phân vị thứ ba của mẫu số liệu ghép nhóm là  \[Q_{3}=30+\dfrac{\dfrac{3\cdot18}{4}-(6+6)}{4}\cdot (35-30)=31{,}875.\]
Khoảng tứ  phân vị của mẫu số liệu ghép nhóm là $\Delta_{Q}=Q_{3}-Q_{1}=8{,}125$.}
\end{ex}

% nguồn: Lop_12_Dai_so_Chuong_3
\begin{ex}%[2D3H1-3]%[ma:d33fdf10]
	Biểu đồ dưới đây biểu diễn số lượng khách hàng đặt bàn qua hình thức trực tuyến mỗi ngày trong quý III năm $2022$ của một nhà hàng. Cột thứ nhất biểu diễn số ngày có từ $1$ đến dưới $6$ lượt đặt bàn, cột thứ hai biễu diễn  số ngày có từ $6$ đến dưới $11$ lượt đặt bàn;$\ldots$
	\begin{center}
		\begin{tikzpicture}[yscale=0.25,>=stealth,line join=round,line cap=round,font=\footnotesize,scale=0.85,line width=1pt]
				\draw[->]  (0,0)--(0,35) node[left]{(\text{Số ngày})};
				\draw[shift={(0,5)}] (0,0)--(-2pt,0) node[left]{\scriptsize  $5$};
				\draw[shift={(0,10)}] (0,0)--(-2pt,0) node[left]{\scriptsize  $10$};
				\draw[shift={(0,15)}] (0,0)--(-2pt,0) node[left]{\scriptsize $15$}; 
				\draw[shift={(0,20)}] (0,0)--(-2pt,0) node[left]{\scriptsize  $20$};
				\draw[shift={(0,25)}]  (0,0)--(-2pt,0) node[left]{\scriptsize $25$};
				 \draw[shift={(0,30)}] (0,0)--(-2pt,0) node[left]{\scriptsize $30$};
				\path (5,34)  node {
				\normalsize{\textbf{Cho mẫu số liệu như biểu đồ  sau}}
					 };
				\draw[dashed,thin,line width=0.01pt] (0,5)--(11,5); 
				\draw[dashed,thin,line  width=0.01pt] (0,10)--(11,10);
				\draw[dashed,thin,line width=0.01pt] (0,15)--(11,15); 
				\draw[dashed,thin,line  width=0.01pt] (0,20)--(11,20);
				 \draw[dashed,thin,line width=0.01pt] (0,25)--(11,25);
				 \draw[dashed,thin,line width=0.01pt] (0,30)--(11,30);
				\draw[line cap=round,pattern=dots] (1,0)--(1,14)--(3,14)  node[midway,above]{$ $}--(3,0)--cycle;
				\draw[line cap=round,pattern=dots] (3,0)--(3,30)--(5,30) node[midway,above]{$  $}--(5,0)--cycle;
				 \draw[line cap=round,pattern=dots] (5,0)--(5,25)--(7,25) node[midway,above]{$ $}--(7,0)--cycle;
				\draw[line cap=round,pattern=dots]  (7,0)--(7,18)--(9,18) node[midway,above]{$ $}--(9,0)--cycle;
				\draw[line  cap=round,pattern=dots] (9,0)--(9,5)--(11,5) node[midway,above]{$ $}--(11,0)--cycle;
				\node [below]  at (2,0){$ [1; 6)$};
				 \node [above] at (2,14){$14$};
				\node [below] at (4,0){$ [6;  11)$};
				\node [above]  at (4,30){$30$};
				\node [below] at  (6,0){$ [11; 16)$};
				 \node [above] at (6,25){$25$};
				\node  [below] at (8,0){$ [16; 21)$};
				\node [above]  at (8,18){$18$};
				 \node [below] at (10,0){$ [21; 26)$};
				 \node [above] at (10,5){$5$};
				\draw[->] (0,0) node [below left=-2pt]{$O$}--(13,0)  node[below]{(\text{Số lượt đặt bàn})};
			\end{tikzpicture}
	\end{center}
	Khoảng tứ phân vị của mẫu số liệu ghép nhóm  cho bởi biểu đồ trên là?
	 \choice
		 {$9{,}5$}
		 {\True $8{,}5$}
		 {$10{,}5$}
		 {$7{,}5$}
	\loigiai{ 
		Dựa vào biểu đồ ta lập được bảng ghép nhóm như sau 
		\begin{center}
			\begin{tabular}{|c|c|c|c|c|c|}
				\hline Lượt đặt bàn & {$[1; 6)$} & {$[6; 11)$}  & {$[11; 16)$} & {$[16; 21)$} & {$[21; 26)$} \\
				\hline Số ngày & $14$ & $30$ & $25$ & $18$ & $5$ \\ 
				 \hline
			\end{tabular}
		\end{center}
		Cỡ mẫu $n=92$ và gọi $x_1, x_2, \ldots, x_{92}$ là mẫu số liệu đã cho 
		 Ta có
		$$ 
		\begin{aligned}
			& x_1, \ldots ., x_{14} \in[1; 6); x_{15}, \ldots,  x_{34} \in[6; 11); x_{35}, \ldots, x_{69} \in[11; 16); \\
			& x_{70}, \ldots, x_{87}  \in[16; 21); x_{88}, \ldots, x_{92} \in[21; 26) .
		\end{aligned}
		 $$
		Tứ phân vị thứ nhất của mẫu số liệu là $\dfrac{x_{23}+x_{24}}{2} \in[6; 11)$. Do đó tứ phân vị thứ nhất của mẫu  số liệu là $Q_1=6+\dfrac{\dfrac{92}{4}-14}{30}(11-6)=7{,}5$.\\
		Tứ phân vị thứ ba của mẫu số liệu là  $\dfrac{x_{69}+x_{70}}{2}$ mà $x_{69} \in[11; 16)$ và $x_{70} \in[16; 21)$. Do đó tứ phân vị thứ ba của mẫu số liệu là $Q_3=16$.\\
		Vậy khoảng tứ phân vị của mẫu  số liệu là $\Delta_Q=Q_3-Q_1=16-7{,}5=8{,}5$.
	} 
\end{ex}

% nguồn: Lop_12_Dai_so_Chuong_3
\begin{ex}%[2D3H1-3]%[ma:60c102e5]
	Cô Minh Hiền rất thích  nhảy hiện đại. Thời gian tập nhảy mỗi ngày trong thời gian gần đây của Cô Minh Hiền được thống kê lại ở bảng sau:
	\begin{center}
		\begin{tabular}{|c|c|c|c|c|c|}
			 \hline
			Thời gian (phút) &$[20;25)$ &$[25;30)$ &$[30;35)$ &$[35;40)$&$[40;45)$ \\ 
			 \hline
			Số ngày &$6$ &$6$  &$4$ &$1$&$1$ \\
			\hline 
		\end{tabular}
	\end{center}
	Khoảng tứ phân vị của mẫu số liệu  ghép nhóm là
	 \choice
		 {$23{,}75$}
		 {$27{,}5$}
		{$31{,}88$} 
		{\True  $8{,}125$}
	 \loigiai{
		Cỡ  mẫu $n=6+6+4+1+1=18$.\\
		Gọi $x_1$; \ldots; $x_{18}$ là mẫu số liệu gốc về thời gian tập nhảy mỗi ngày của Cô Minh Hiền được xếp theo thứ tự  không giảm.\\
		 Ta có
		$x_1;\ldots;x_6\in [20; 25)$, $x_7;\ldots;x_{12}\in [25; 30)$,  $x_{13};\ldots;x_{16}\in [30; 35)$,
		$x_{17}\in [35; 40)$, $x_{18}\in [40; 45)$. 
		\begin{itemize}
			\item Tứ phân  vị thứ nhất của mẫu số liệu gốc là $x_5\in [20;25)$.\\
			Do đó, tứ phân vị thứ nhất của mẫu số liệu  ghép nhóm là $$Q_1=20+\dfrac{\dfrac{18}{4}}{6}(25-20)=23{,}75.$$
			\item Tứ phân vị thứ ba của  mẫu số liệu gốc là $x_{14}\in [30;35)$.\\
			Do đó, tứ phân  vị thứ ba của mẫu số liệu ghép nhóm là
			$$Q_3=30+\dfrac{\dfrac{3\cdot  18}{4}-(6+6)}{4}(35-30)=31{,}875.$$
			Khoảng tứ phân vị của mẫu  số liệu ghép nhóm là
			 $$\Delta_Q=Q_3-Q_1=31{,}875-23{,}75=8{,}125.$$
		\end{itemize}
	 }
\end{ex}

% nguồn: Lop_12_Dai_so_Chuong_3
\begin{ex}%[2D3H1-3]%[ma:f12b17ad]
	Dũng là học sinh rất giỏi chơi rubik, bạn có thể  giải nhiều loại khối rubik khác nhau. Trong một lần tập luyện giải khối rubik $3 \times 3$, bạn Dũng đã tự thống kê lại thời gian giải rubik trong $25$ lần giải liên tiếp ở bảng sau
	\begin{center}
		\begin{tabular}{|c|c|c|c|c|c|}
			\hline \begin{tabular}{c}
			Thời gian giải  \\
			 rubik (giây)
		\end{tabular} & {$[8; 10)$} & {$[10; 12)$} & {$[12; 14)$} & {$[14; 16)$} & {$[16; 18)$} \\
		\hline Số lần & $4$ & $6$ & $8$ &  $4$ & $3$ \\
		\hline 
	\end{tabular}
\end{center}
Khoảng tứ phân vị của mẫu số liệu ghép nhóm  là
 \choice
	{$10{,}75$} 
	{$1{,}75$} 
	{\True  $3{,}63$}
	 {$14{,}38$}
\loigiai{ 
	 Cỡ mẫu $n=25$.\\
	Gọi $x_{1} $; $x_{2} $; $\ldots $; $x_{25}$ là mẫu số liệu gốc về thời gian giải rubik trong $25$ lần của bạn Dũng được  xếp theo thứ tự không giảm.\\
	Ta có $x_{1} , \ldots , x_{4} \in[8; 10)  $; $x_{5} , \ldots , x_{10} \in[10; 12) $; $x_{11} , \ldots , x_{18} \in[12; 14) $; $x_{19} , \ldots , x_{22} \in[14; 16)$; $x_{23} , \ldots , x_{25} \in[16; 18)$.\\
	 Tứ phân vị thứ nhất của mẫu số liệu gốc là $\dfrac{1}{2}(x_{6}+x_{7}) \in[10; 12)$. Do đó, tứ phân vị thứ nhất của mẫu số liệu ghép nhóm là \[Q_{1}=10+\dfrac{\dfrac{25}{4}-4}{6}\cdot(12-10)=10{,}75.\]
	Tứ phân vị thứ ba của mẫu  số liệu gốc là $x_{19} \in[14; 16)$.
	Do đó, tứ phân vị thứ  ba của mẫu số liệu ghép nhóm là
	\[Q_{3}=14+\dfrac{\dfrac{3\cdot  25}{4}-(4+6+8)}{4}\cdot (16-14)=14{,}375.\]
Khoảng tứ phân vị của mẫu số  liệu ghép nhóm là $\Delta_{Q}=Q_{3}-Q_{1}=3{,}63$.}
\end{ex}


%%%%%%%%%%%%%%%% ĐÚNG/SAI %%%%%%%%%%%%%%%%

% nguồn: Lop_12_Dai_so_Chuong_3
\begin{ex}%[2D3H1-3]%[ma:f4ed13b4]
	Phòng quản lí đào tạo trường Đại học Kinh tế quốc dân thống kê số giờ làm thêm của một nhóm sinh viên  năm thứ tư của trường thu được kết quả như bảng sau:
	\begin{center}
		\resizebox{0.8\linewidth}{!}{\begin{tabular}{|c|c|c|c|c|c|}
		 \hline
		\thead{\textbf{Số giờ làm thêm (giờ/tuần)}} & $[9; 12)$ & $[12; 15)$ &  $[15; 18)$ & $[18; 21)$ & [$21; 24)$ \\
		\hline 
		\thead{\textbf{Số  sinh viên}} & $6$ & $12$ & $4$ & $2$ & $1$ \\
		\hline 
	\end{tabular}}
\end{center}
 \choiceTF
	{Số giờ làm thêm trung bình của nhóm sinh viên trên trong một tuần  là~$16{,}5$ giờ}
	{\True Giá trị đại diện  của nhóm $[9;12)$ là~$10{,}5$}
	{Tứ phân vị thứ  ba là~$15{,}65$}
	{Nhóm chứa trung vị  là~$[15;18)$}
 \loigiai{
	\begin{itemchoice}
		\itemch Cỡ mẫu:  $n=6+12+4+2+1=25$. \\
		Ta có  bảng sau:
		\begin{center}
			\resizebox{0.8\linewidth}{!}{\begin{tabular}{|c|c|c|c|c|c|}
			\hline 
			\thead{\textbf{Số giờ làm thêm (giờ/tuần)}} & $[9; 12)$  & $[12; 15)$ & $[15; 18)$ & $[18; 21)$ & $[21; 24)$ \\
			 \hline
			\thead{\textbf{Giá trị đại diện}} & $10{,}5$ & $13{,}5$  & $16{,}5$ & $19{,}5$ & $22{,}5$ \\
			\hline 
			\thead{\textbf{Số sinh viên}} & $6$ & $12$  & $4$ & $2$ & $1$ \\
			 \hline
		\end{tabular}}
	\end{center}
	Số giờ làm thêm trung bình của nhóm sinh viên trên trong một tuần  là
	$$\overline{x}=\dfrac{6\cdot 10{,}5+12\cdot  13{,}5+4\cdot 16{,}5+2\cdot 19{,}5+1\cdot 22{,}5}{25}=14{,}1 \text{ (giờ).}$$
	\itemch Giá  trị đại diện của nhóm $\left[9;12\right)$ là~$10{,}5$.
	\itemch Giả sử $x_1, x_2,\ldots, x_{25}$ số  giờ làm thêm của các sinh viên trong mẫu số liệu trên và dãy này đã được sắp xếp theo thứ tự không giảm. \\
	Khi đó, trung vị của mẫu số liệu là~$x_{13}$ và tứ phân vị thứ ba là~$\dfrac{1}{2} \left(x_{19}+x_{20} \right)$. Vì $x_{19}, x_{20}$ đều thuộc nhóm $\left[15;18\right)$ nên nhóm này chứa tứ phân vị thứ ba. Do  đó, tứ phân vị thứ ba là:
	\[Q_3=15+\dfrac{\dfrac{3\cdot  25}{4}-12-6}{4} \cdot \left(18-15\right)=15{,}5625.
	\] 
	\itemch Vì $x_{13}$ thuộc  nhóm $\left[12;15\right)$ nên nhóm chứa trung vị là~nhóm $\left[12;15\right)$.
	\end{itemchoice}
 }
\end{ex}

% nguồn: Lop_12_Dai_so_Chuong_3
\begin{ex}%[2D3H1-3]%[ma:fdd39a9c]
	Biểu đồ dưới đây thống kê thời gian tập thể dục buổi sáng mỗi ngày trong tháng $9 / 2022$  của bác Bình và bác An.
	\begin{center}
		\begin{tikzpicture}[>=stealth,declare function={ox=-1.5; kc=2.0; rc=1.0; kcy=1;}]
				\foreach[count=\i]  \yt in {5,10,15,20,25,30}{
					\draw[gray!35]  (0,{\i*kcy}) node[left,text=black]{$\yt$}--+(10,0);}
				\foreach[count=\i] \yt/\gt/\yh/\gh/\xt in  {1/{5}/0.02/{0}/{$[15; 20)$},2.1/{12}/5/{25}/{$[20; 25)$},1.8/{8}/1/{5}/{$[25; 30)$},0.6/{3}/0.02/{0}/{$[30; 35)$},0.2/{2}/0.02/{0}/{$[35; 40)$}}{
					\draw[blue,fill=blue] ({ox+\i*kc},0)  rectangle +(rc,{\yt*kcy});
					 \path ({ox+\i*kc+0.5*rc},{\yt*kcy}) node[above]{$\gt$};
					\draw[orange,fill=orange] ({ox+\i*kc+rc},0)  rectangle +(rc,{\yh*kcy});
					\path  ({ox+\i*kc+1.5*rc},{\yh*kcy}) node[above]{$\gh$};
					\path  ({ox+\i*kc+rc},0) node[below]{\xt};
				} 
				\draw[->] (0,0)--(11.5,0) node[shift={(-6pt,-6pt)}]{Thời gian}; 
				 \draw (11.5,-0.5) node[shift={(-6pt,-6pt)}]{(phút)};
				\draw[->] (0,0)--(0,7) node[left,anchor=east]{Số ngày}; 
				\begin{scope}[shift={(10,5)}]
					\draw[fill=blue] (0,0) rectangle  +(0.3,0.3);
					 \node[right] at (0.6,0.25) {Bác Bình};
					\draw[fill=orange] (0,-0.8) rectangle +(0.3,0.3); 
					\node[right] at (0.6,-0.55) {Bác  An};
				\end{scope}
			\end{tikzpicture}
	\end{center}
	 Khi đó
	 \choiceTF
		{\True Khoảng biến thiên của  mẫu số liệu ghép nhóm về thời gian tập thể dục buổi sáng của bác Bình là $25$ (phút)}
		{Khoảng tứ phân vị của mẫu số liệu ghép nhóm về thời gian  tập thể dục buổi sáng của bác An là $\Delta_{Q}=2$}
		{\True Tứ phân vị thứ ba của mẫu số liệu ghép nhóm về thời gian tập thể dục buổi sáng của bác Bình là $Q_{3}^{\prime}=\dfrac{455}{16}$} 
		{Khoảng tứ phân vị của mẫu số liệu  ghép nhóm về thời gian tập thể dục buổi sáng mỗi ngày của bác An lớn hơn bác Bình}
	 \loigiai{
		\begin{itemchoice}
			\itemch  Ta có bảng sau
			\begin{center}
				\begin{tabular}{|c|c|c|c|c|c|}
					\hline \begin{tabular}{c}
					Thời  gian \\
					 (phút)
				\end{tabular} &$[15; 20)$ &$[20; 25)$ &$[25; 30)$ &$[30; 35)$ &$[35; 40)$ \\
				\hline Bác  Bình & $5$ & $12$ & $8$ & $3$ & $2$ \\
				\hline Bác An & $0$ & $25$ & $5$ & $0$ &$  0$ \\
				 \hline
			\end{tabular}
		\end{center}
		 \itemch Cỡ mẫu $n=30$;\\
		Gọi $x_{1}; x_{2}; \ldots; x_{30}$ là mẫu số liệu gốc về thời gian tập thể dục buổi sáng mỗi ngày  của bác An được xếp theo thứ tự không giảm.
		 Ta có $x_{1} $, $x_{2} $, $\ldots$ , $x_{25} \in[20; 25) $; $x_{26} $, $\ldots $, $x_{30} \in[25; 30)$.
		Tứ phân vị thứ nhất của mẫu số liệu gốc là $x_{8}  \in[20; 25)$. \\
		 Do đó, tứ phân vị thứ nhất của mẫu số liệu ghép nhóm là $Q_{1}=20+\dfrac{\tfrac{30}{4}}{25}\cdot(25-20)=\dfrac{43}{2}$.\\
		Tứ phân vị thứ ba của mẫu số  liệu gốc là $x_{23} \in[20; 25)$. \\
		Do đó, tứ phân vị thứ ba của mẫu số liệu ghép nhóm là $Q_{3}=20+\dfrac{\dfrac{3\cdot  30}{4}}{25}(25-20)=\dfrac{49}{2}$.\\
		Khoảng  tứ phân vị của mẫu số liệu ghép nhóm là $\Delta_{Q}=Q_{3}-Q_{1}=3$.\\
		Gọi $y_{1} $; $y_{2} $; $\ldots $; $y_{30}$  là mẫu số liệu gốc về thời gian tập thể dục buổi sáng mỗi ngày của bác Bình được xếp theo thứ tự không giảm.\\
		Ta  có
		 $y_{1} ,\, y_{2} ,\, \ldots ,\, y_{5} \in[15; 20) $; $y_{6} ,\, \ldots ,\, y_{17} \in[20; 25) $; $y_{18} ,\, \ldots ,\, y_{25} \in[25; 30) $; \\
		$y_{26}  ,\, y_{27} ,\, y_{28} \in[30; 35)$; $y_{29} ,\, y_{30} \in[35; 40)$.\\
		Tứ phân vị thứ nhất của mẫu  số liệu gốc là $y_{8} \in[20; 25)$. Do đó, tứ phân vị thứ nhất của mẫu số liệu ghép nhóm là $Q_1'=20+\dfrac{\dfrac{30}{4}}{12}(25-20)=\dfrac{185}{8}$.
		\itemch Tứ phân vị thứ ba của mẫu số liệu gốc là $y_{23} \in[25; 30)$.  Do đó, tứ phân vị thứ ba của mẫu số liệu ghép nhóm là $Q_3^{\prime}=25+\dfrac{\dfrac{3.30}{4}-(5+12)}{8}(30-25)=\dfrac{455}{16}$.
		\itemch Vậy khoảng tứ phân vị của mẫu số liệu ghép nhóm về thời gian tập thể dục buổi sáng mỗi ngày của bác Bình lớn hơn bác An. 
	\end{itemchoice}
 }
\end{ex}

% nguồn: Lop_12_Dai_so_Chuong_3
\begin{ex}%[2D3H1-3]%[ma:d5e9ceec]
	Bảng sau biểu diễn mẫu số liệu ghép nhóm vể độ tuổi của cư dân trong một  khu phố.\\
	\begin{center}
		\begin{tabular}{|c|c|}
			\hline Nhóm &  Tần số \\
			\hline$[20; 30)$  & $25$ \\
			\hline$[30; 40)$ & $20$  \\
			\hline$[40;  50)$ & $20$ \\
			\hline$[50; 60)$ & $15$  \\
			\hline$[60; 70)$ & $14$ \\ 
			\hline$[70; 80)$ & $6$ \\ 
			\hline  & $n=100$ \\
			\hline 
		\end{tabular}
	\end{center}
	Xét  tính đúng sai của các khẳng định sau
	\choiceTF 
		{\True Khoảng biến thiên của mẫu số liệu ghép nhóm  trên là $R=60$}
		{\True Tứ phân  vị thứ nhất là $Q_{1}=35$}
		{Tứ  phân vị thứ ba là $Q_{3}=\dfrac{160}{3}$}
		{\True Khoảng tứ phân vị của mẫu số liệu ghép  nhóm trên là $\Delta_{Q}=\dfrac{65}{3}$}
	 \loigiai{
		\begin{itemchoice}
			\itemch Khoảng biến thiên của mẫu  số liệu ghép nhóm trên là $R=80-20=60$
			\itemch Số phần tử của  mẫu là $n=100$.\\
			Tần số tích lũy của các nhóm  lần lượt là
			 \[c f_{1}=25,\, c f_{2}=45,\,c f_{3}=65,\,c f_{4}=80,\,c f_{5}=94,\,c f_{6}=100.\]
			 Ta có $\dfrac{n}{4}=\dfrac{100}{4}=25$ suy ra nhóm $1$ là nhóm đầu tiên có tần số tích lũy lớn hơn hoặc bằng $25$. Xét nhóm $1$ là nhóm $[20;30]$ có $s=20$, $h=10$, $n_{1}=25$.\\
			Ta có tứ phân vị thứ nhất là $Q_{1}=s+\left(\dfrac{25-c f_{0}}{n_{1}}\right) \cdot h=25+\left(\dfrac{25-0}{25}\right) \cdot 10=35$. 
			\itemch Ta có $\dfrac{3 n}{4}=\dfrac{3\cdot100}{4}=75$ mà $65<75<80$ suy ra nhóm $4$  là nhóm đầu tiên có tần số tích lũy lớn hơn hoặc bẳng $80$. \\
			Xét nhóm $4$ là nhóm $[50; 60]$ có $t=50$,  $l=10$, $n_{4}=15$ và nhóm $3$ là nhóm $[40; 50]$ có $c f_{3}=65$.\\
			 Ta có tứ phân vị thứ ba là $Q_{3}=t+\left(\dfrac{75-c f_{3}}{n_{4}}\right) \cdot l=50+\left(\dfrac{75-65}{15}\right) \cdot 10=\dfrac{170}{3}$.
			\itemch Khoảng tứ phân vị của mẫu số liệu ghép nhóm trên là  $Q_{3}-Q_{1}=\dfrac{170}{3}-35=\dfrac{65}{3}$.\\
	\end{itemchoice}}
\end{ex}

% nguồn: Lop_12_Dai_so_Chuong_3
\begin{ex}%[2D3H1-3]%[2D3H2-2]%[ma:697a60b1]
	Bảng sau thống kê lại tổng số giờ nắng trong tháng 6 của các năm từ 2002 đến 2021 tại  hai trạm quan trắc đặt ở Nha Trang và Quy Nhơn.\\
	\centerline{ 
	\begin{tabular}{|c|c|c|c|c|c|c|}\hline
		Số giờ nắng& $[130;160)$& $[160;190)$& $[190;220)$&  $[220;250)$& $[250;280)$& $[280;310)$\\\hline
		Số năm ở Nha  Trang& $1$& $1$& $1$& $8$& $7$& $2$\\
		Số năm ở Quy Nhơn& $0$& $1$& $2$& $4$&  $10$& $3$\\\hline
	\end{tabular}
	}\\ 
	Xét tính  đúng sai của các khẳng định sau:
	 \choiceTF
		{Xét số liệu ở Nha Trang thì khoảng tứ phân vị của mẫu số  liệu ghép nhóm là: $32{,}64$}
		{\True Nếu so sánh  theo khoảng tứ phân vị thì số giờ nắng trong tháng 6 của Quy Nhơn đồng đều hơn}
		{\True Xét số  liệu của Quy Nhơn ta có độ lệch chuẩn của mẫu số liệu ghép nhóm (làm tròn kết quả đến hàng phần trăm) là: $30{,}59$}
		{Nếu so sánh theo độ lệch chuẩn thì số giờ nắng trong tháng 6 của  Nha Trang đồng đều hơn}
	 \loigiai{
		\begin{itemchoice}
			\itemch Sai:  Cỡ mẫu: $n=20$.\\
			Gọi $x_1;x_2;\ldots;x_{20}$ là mẫu số liệu gốc về số giờ nắng trong tháng 6 trong 20 năm của Nha Trang được xếp theo  thứ tự không giảm.\\
			Ta có: $x_1\in  [130;160);x_2\in [160;190);x_3\in [190;220);x_4;\ldots;x_{11}\in [220;250);x_{12};\ldots;$.\\
			$x_{18}\in  [250;280){;}x_{19};x_{20}\in [280;310)$.\\
			Tứ phân vị thứ  nhất của mẫu số liệu gốc là $\dfrac{1}{2}(x_5+x_6)\in [220;250)$. \\
			Do đó, tứ phân vị  thứ nhất của mẫu số liệu ghép nhóm là: \\
			 $\mathrm{Q}_1=220+\dfrac{\dfrac{20}{4}-(1+1+1)}{8}(250-220)=227{,}5$.\\
			Tứ phân vị thứ ba của mẫu số liệu gốc là $\dfrac{1}{2}(x_{15}+x_{16})\in  [250;280)$. \\
			 Do đó, tứ phân vị thứ ba của mẫu số liệu ghép nhóm là: \\
			$\mathrm{Q}_3=250+\dfrac{\dfrac{3\cdot  20}{4}-(1+1+1+8)}{7}(280-250) =\dfrac{1870}{7}$.\\
			Khoảng tứ  phân vị của mẫu số liệu ghép nhóm là: $\Delta_{\mathrm{Q}}=\mathrm{Q}_3-\mathrm{Q}_1=39{,}64$.
			\itemch Đúng: Gọi $y_1; y_2;\ldots; y_{50}$ là mẫu số liệu  gốc về số giờ nắng trong tháng 6 trong 20 năm của Quy Nhơn được xếp theo thứ tự không giảm.\\
			 Ta có:\\
			$y_1;\in [160;190); y_2;y_3\in [190;220); y_4;\ldots; y_7\in [220;250);  y_8;\ldots; y_{17}\in [250;280);$\\
			 $y_{4=18};\ldots; y_{20}\in [280;310)$.\\
			Tứ  phân vị thứ nhất của mẫu số liệu gốc là $\dfrac{1}{2}(y_5+y_6)\in [220;250)$. \\
			Do đó, tứ phân vị  thứ nhất của mẫu số liệu ghép nhóm là: \\
			$Q'_1=220+\dfrac{\dfrac{20}{4}-(1+2)}{4}(250-220)=235$.\\ 
			Tứ phân vị thứ  ba của mẫu số liệu gốc là $\dfrac{1}{2}(y_{15}+y_{16})\in [250;280)$. \\
			Do đó, tứ phân  vị thứ ba của mẫu số liệu ghép nhóm là: \\
			$Q'_3=250+\dfrac{\dfrac{3\cdot  20}{4}-(1+2+4)}{10}(280-250) =274$.\\
			Khoảng tứ phân vị của mẫu số  liệu ghép nhóm là: $\Delta'_Q=Q'_3-Q'_1=39$.\\
			Vậy nếu so sánh theo  khoảng tứ phân vị thì số giờ nắng trong tháng 6 của Quy Nhơn đồng đều hơn.\\
			Xét số liệu của Nha Trang:\\ 
			 Số trung bình: $\overline{x}_{_X}=\dfrac{1\cdot 145+1\cdot 175+1\cdot 205+8\cdot 235+7\cdot 265+2\cdot 295}{20}=242{,}5$.\\
			Độ lệch  chuẩn: $\sigma_X=\sqrt{\dfrac{1\cdot 145^2+1\cdot 175^2+1\cdot 205^2+8\cdot 235^2+7\cdot 265^2+2\cdot 295^2}{20}-242{,}5^2} \approx 35{,}34$.
			 \itemch Đúng: Xét số liệu của Quy Nhơn:\\
			Số trung bình: $\overline{x}_{_Y}=\dfrac{1\cdot 175+2\cdot  205+4\cdot 235+10\cdot 265+3\cdot 295}{20}=253$.\\
			Độ lệch chuẩn: $\sigma_Y=\sqrt{\dfrac{1\cdot 175^2+2\cdot 205^2+4\cdot 235^2+10\cdot 265^2+3\cdot 295^2}{20}-253^2} \approx  30{,}59$.
			\itemch Sai: Vậy nếu so sánh theo độ lệch chuẩn thì số giờ nắng trong  tháng 6 của Quy Nhơn đồng đều hơn.
		\end{itemchoice}
	} 
\end{ex}

% nguồn: Lop_12_Dai_so_Chuong_3
\begin{ex}%[2D3H1-3]%[ma:c65fe323]
	Bảng dưới đây biểu diễn mẫu số liệu ghép nhóm về nhiệt độ không  khí trung bình các tháng năm 2021 tại Hà Nội và Huế (đơn vị: độ C).\\
	 \centerline{
	\begin{tabular}{|>{\centering\arraybackslash}m{25mm}|>{\centering\arraybackslash}m{15mm}|>{\centering\arraybackslash}m{15mm}|}\hline
		Nhóm& Hà Nội& Huế\\\hline 
		$[16{,}8;19{,}8)$& $2$&  $1$\\
		$[19{,}8;22{,}8)$& $3$& $2$\\ 
		 $[22{,}8;25{,}8)$& $2$& $3$\\
		$[25{,}8;28{,}8)$& $1$& $2$\\ 
		$[28{,}8;31{,}8)$&  $4$& $4$\\\hline
		&$n=12$& $n=12$\\\hline 
	\end{tabular}
	} 
	Xét tính đúng sai của các khẳng định  sau:
	 \choiceTF
		{\True Khoảng tứ phân vị của mẫu số liệu ghép nhóm trên  số liệu ở Hà Nội là: $8{,}75$}
		{Độ lệch chuẩn của mẫu số liệu ghép  nhóm trên số liệu ở Hà Nội (làm tròn kết quả đến hàng phần trăm) là: $3{,}56$}
		{Phương sai của mẫu số liệu ghép nhóm trên số liệu  ở Huế là: $12{,}75$}
		{\True Huế  có nhiệt độ không khí trung bình tháng đồng đều hơn vì độ lệch chuẩn nhỏ hơn}
	 \loigiai{
		 \centerline{
		\begin{tabular}{|>{\centering\arraybackslash}m{25mm}|>{\centering\arraybackslash}m{15mm}|>{\centering\arraybackslash}m{15mm}|>{\centering\arraybackslash}m{15mm}|>{\centering\arraybackslash}m{15mm}|}\hline
			Nhóm& Hà Nội& Tần  số tích luỹ& Huế& Tần số tích luỹ\\\hline
			$[16{,}8;19{,}8)$& $2$&  $2$& $1$& $1$\\
			 $[19{,}8;22{,}8)$& $3$& $5$& $2$& $3$\\
			$[22{,}8;25{,}8)$& $2$& $7$&  $3$& $6$\\
			 $[25{,}8;28{,}8)$& $1$& $8$& $2$& $8$\\
			$[28{,}8;31{,}8)$&  $4$& $12$& $4$& $12$\\\hline
			&$n=12$&& $n=12$&\\\hline 
		\end{tabular}
		 }\\
		\begin{itemchoice}
			\itemch Đúng:  Xét số liệu ở Hà Nội:\\
			Khoảng biến  thiên: $R=31{,}8-16{,}8=15$.\\
			Số phần tử  của mẫu là $n=12$.\\
			Tần số tích lũy của các nhóm lần lượt  là $cf_1=2,cf_2=5,cf_3=7,cf_4=8,cf_5=12$.\\
			Ta có: $\dfrac{n}{4}=\dfrac{12}{4}=3$ mà $2<3<5$ suy ra nhóm 2 là nhóm đầu tiên có tần số tích lũy lớn  hơn hoặc bằng 3. \\
			Xét  nhóm 2 là nhóm $[19{,}8;22{,}8)$ có $s=19{,}8$, $h=3$, $n_2=3$ và nhóm 1 là nhóm $[16{,}8;19{,}8)$ có $cf_1=2$.\\
			Ta có tứ phân vị thứ  nhất là: $\mathrm{Q}_1=s+\left(\dfrac{3-cf_1}{n_2}\right)\cdot h =19{,}8+\left(\dfrac{3-2}{3}\right)\cdot 3=20{,}8$.\\
			Ta có: $\dfrac{3n}{4}=\dfrac{3\cdot 12}{4}=9$ mà $8<9<12$ suy ra nhóm 5 là nhóm đầu tiên có tần số tích  lũy lớn hơn hoặc bằng 9. \\
			Xét nhóm 5  là nhóm $[28{,}8;31{,}8)$ có $t=28{,}8$, $l=3$, $n_5=4$ và nhóm 4 là nhóm $[25{,}8;28{,}8)$ có $cf_4=8$.\\
			 Ta có tứ phân vị thứ ba là: $\mathrm{Q}_3=t+\left(\dfrac{9-cf_4}{n_5}\right)\cdot l =28{,}8+\left(\dfrac{9-8}{4}\right)\cdot 3=29{,}55$.\\
			Khoảng tứ phân vị của mẫu số  liệu ghép nhóm trên là: $\mathrm{Q}_3-\mathrm{Q}_1=29{,}55-20{,}8=8{,}75$.
			\itemch Sai:  Số trung bình cộng của mẫu số liệu ghép nhóm là: \\
			$\overline{x}_1=\dfrac{2\cdot 18{,}3+3\cdot  21{,}3+2\cdot 24{,}3+27{,}3+4\cdot 30{,}3}{12}=24{,}8$.\\
			Phương sai của mẫu số liệu ghép  nhóm là:\\
			 $s_1^2=\dfrac{2{(18{,}3-24{,}8)^2}+3(21{,}3-24{,}8)^2+2(24{,}3-24{,}8)^2+(27{,}3-24{,}8)^2+4(30{,}3-24{,}8)^2}{12}=20{,}75$.\\
			Độ lệch chuẩn của mẫu số liệu  ghép nhóm là: $s_1=\sqrt{s_1^2}=\sqrt{20{,}75}\approx 4{,}56$.
			\itemch Sai: Xét số liệu ở  Huế:\\
			 Khoảng biến thiên: $R=31{,}8-16{,}8=15$.\\
			Số phần tứ của  mẫu là $n=12$.\\
			Tần số tích lũy của các nhóm lần lượt là $cf_1=1$,  $cf_2=3$, $cf_3=6$, $cf_4=8$, $cf_5=12$.\\
			Ta có: $\dfrac{n}{4}=\dfrac{12}{4}=3$ suy ra nhóm 2 là nhóm  đầu tiên có tần số tích lũy lớn hơn hoặc bằng 3. \\
			Xét nhóm 2 là nhóm $[19{,}8;22{,}8)$ có $s=19{,}8$, $h=3$, $n_2=2$ và nhóm 1 là nhóm  $[16{,}8;19{,}8)$ có $cf_1=1$.\\
			Ta có tứ phân vị thứ  nhất là: $\mathrm{Q}_1=s+\left(\dfrac{3-cf_1}{n_2}\right)\cdot h =19{,}8+\left(\dfrac{3-1}{2}\right)\cdot 3 =22{,}8$.\\
			Ta có: $\dfrac{3n}{4}=\dfrac{3\cdot  12}{4}=9$ mà $8<9<12$ suy ra nhóm 5 là nhóm đầu tiên có tần số tích lũy lớn hơn hoặc bằng 9. \\
			Xét  nhóm 5 là nhóm $[28{,}8;31{,}8)$ có $t=28{,}8$, $l=3$, $n_5=4$ và nhóm 4 là nhóm $[25{,}8;28{,}8)$ có $cf_4=8$.\\
			Ta có tứ  phân vị thứ ba là: $\mathrm{Q}_3=t+\left(\dfrac{9-cf_4}{n_5}\right)\cdot l =28{,}8+\left(\dfrac{9-8}{4}\right)\cdot 3 =29{,}55$.\\
			Khoảng  tứ phân vị của mẫu số liệu ghép nhóm trên là: $\mathrm{Q}_3-\mathrm{Q}_1=29{,}55-22{,}8=6{,}75$.\\
			Số  trung bình cộng của mẫu số liệu ghép nhóm là: \\
			$\overline{x}_2=\dfrac{18{,}3+2\cdot 21{,}3+3\cdot  24{,}3+2\cdot 27{,}3+4\cdot 30{,}3}{12} =25{,}8$.\\
			Phương sai của mẫu số liệu  ghép nhóm là:\\
			$s_2^2=\dfrac{(18{,}3-25{,}8)^2+3(21{,}3-25{,}8)^2+3(24{,}3-25{,}8)^2+2(27{,}3-25{,}8)^2+4(30{,}3-25{,}8)^2}{12}=15{,}75$. 
			\itemch Đúng: Độ lệch chuẩn của mẫu số liệu ghép  nhóm là: $s_2=\sqrt{s_2^2}=\sqrt{15{,}75}\approx 3{,}97$.\\
			Huế có nhiệt độ không khí trung bình tháng đồng đều hơn  vì độ lệch chuẩn nhỏ hơn.
		\end{itemchoice}
	 }
\end{ex}

% nguồn: Lop_12_Dai_so_Chuong_3
\begin{ex}%[2D3H1-3]%[ma:1657fdae]
	 Bạn Trang thống kê lại chiều cao (đơn vị: cm) của các bạn học sinh nữ lớp $12C$ và lớp $12D$ ở bảng sau\\
	 \resizebox{\linewidth}{!}{
	\begin{tabular}{|c|c|c|c|c|c|c|}
		\hline \diagbox{Lớp}{Chiều cao (cm)} &  $[155;160)$ & $[160;165)$ & $[165;170)$ & $[170;175)$ & $[175;180)$ & $[180;185)$ \\
		\hline $12C$ &  $2$ & $7$ & $12$ & $3$ & $0$ & $1$ \\
		\hline $12D$ & $5$ & $9$ & $8$ & $2$ & $1$  & $0$ \\
		\hline 
	\end{tabular}
	} 
	Xét tính đúng sai của các khẳng định sau: 
	\choiceTF 
		{Chiều cao cao nhất của  các bạn học sinh trong lớp $12D$ là $185$ cm}
		{\True Khoảng biến thiên của chiều cao của các bạn học sinh nữ lớp $12 C$ là  $30$ cm}
		{\True Khoảng biến thiên của chiều cao của các  bạn học sinh nữ lớp $12D$ là $25$ cm}
		{Chiều  cao của học sinh lớp $12 C$ có độ phân tán bé hơn}
	 \loigiai{
		Khoảng biến thiên của chiều cao của các bạn học sinh nữ  lớp $12 C$ là $185-155=30$ (cm).\\
		Khoảng biến thiên của chiều cao của các  bạn học sinh nữ lớp $12D$ là $180-155=25$ (cm).\\
		Vậy nếu so sánh theo khoảng biến thiên thì chiều  cao của học sinh lớp $12 C$ có độ phân tán lớn hơn.
		\begin{itemchoice}
			\itemch Chiều cao cao nhất của các bạn học sinh trong lớp $12D$  là $180$ cm.
			\itemch Khoảng biến thiên của chiều cao của các bạn học sinh nữ lớp $12  C$ là $30$ cm.
			\itemch Khoảng biến thiên của chiều cao của các  bạn học sinh nữ lớp $12D$ là $25$ cm.
			\itemch Chiều cao của  học sinh lớp $12D$ có độ phân tán bé hơn.
	\end{itemchoice}}
\end{ex}

% nguồn: Lop_12_Dai_so_Chuong_3
\begin{ex}%[2D3H1-3]%[ma:45bb0e76]
	Bảng sau thống kê tổng lượng mưa (đơn vị: mm) đo được  vào tháng $7$ từ năm $2\,002$ đến $2\,021$ tại một trạm quan trắc đặt ở Cà Mau
	\begin{center}
		\begin{tabular}{|c|c|c|c|c|c|c|c|c|c|}
			\hline 
			$341{,}4$  & $187{,}1$& $242{,}2$& $5{,22}9$& $251{,}4$& $432{,}2$& $200{,}7$& $388{,}6$& $258{,}4$& $288{,}5$\\
			\hline 
			 $298{,}1$ & $413{,}5$& $413{,}5$& $332$& $421$& $475$& $400$& $305$& $520$& $1477$\\
			\hline 
		\end{tabular}
	\end{center}
	Chia  mẫu số liệu trên thành $4$ nhóm như bảng
	\begin{center}
		\begin{tabular}{|c|c|c|c|c|}
			\hline 
			Lượng mưa & $[140;240)$& $[240;340)$& $[340;440)$& $[440;540)$\\ 
			\hline 
			Số tháng & $3$&  $7$& $7$& $3$\\
			\hline 
		\end{tabular}
	\end{center}
	\choiceTF 
		{\True Khoảng biến  thiên của mẫu số liệu $375,9$\,(mm)}
		{Khoảng tứ phân vị của  mẫu số liệu là $\Delta_Q=\dfrac{1827}{100}$}
		{\True  Khoảng biến thiên của mẫu số liệu ghép nhóm $400$ (mm)}
		{\True  Khoảng tứ phân vị của mẫu số liệu ghép nhóm là $\Delta_Q ^{'}=\dfrac{1000}{7}$}
	\loigiai{ 
		\begin{itemchoice}
			\itemch Sắp  xếp lại mẫu số liệu theo thứ tự tăng dần
			\begin{center}
				$\begin{array}{cccccccccc}
				147&187{,}1&200{,}7 &242{,}2&251{,}4&258{,}4& 288{,}5& 298{,}1&  305&332\\
				341{,}4&  388{,}6& 400& 413{,}5& 413{,}5& 421& 432{,}2&475& 520& 522{,}9
			\end{array}$
		\end{center}
		Khoảng biến thiên của mẫu số liệu $522,9-147=375,9$  (mm).
		\itemch Tứ phân vị thứ nhất của mẫu số liệu  là trung vị của
		\begin{center}
			$\begin{array}{cccccccccc}
			147 &187{,}1 &200{,}7 &242{,}2 &251{,}4 &258{,}4 &288{,}5 &298{,}1 &305 &332 
		\end{array}$
	\end{center}
	nên $Q_1=\dfrac{6276}{25}$.\\ 
	Tứ phân  vị thứ ba của mẫu số liệu là trung vị của
	\begin{center}
		$\begin{array}{cccccccccc}
		341{,}4 &388{,}6 &400 &413{,}5 &413{,}5 &421  &432{,}2 &475 &520 &522{,}9
	\end{array}$
	\end{center}
	 nên $Q_3=\dfrac{43281}{100}$.\\
	 Khoảng tứ phân vị của mẫu số liệu là $\Delta_Q=Q_3-Q_1=\dfrac{18177}{100}$.
	\itemch Khoảng biến thiên của mẫu số liệu ghép nhóm $540-140=400$  (mm).
	Cỡ mẫu  $n=20$;\\
	Gọi $x_1;x_2;\ldots;x_{20}$ là  mẫu số liệu gốc về lượng mưa đo được vào tháng $7$ từ năm $2\,002$ đến $2\,021$ tại một trạm quan trắc đặt ở Cà Mau được xếp theo thứ tự không giảm.\\
	 Ta có
	\begin{itemize}
		\item $x_1;\ldots;x_3  \in [140;240)$;
		\item $x_4;\ldots;x_{10} \in  [240;340)$;
		\item $x_{11};\ldots;x_{17} \in  [340;440)$;
		\item  $x_{18};\ldots;x_{20} \in [440;540)$.
	\end{itemize}
	Tứ phân vị thứ nhất của  mẫu số liệu gốc là $\dfrac{1}{2} \left(x_5+x_6 \right)\in [240;340)$.\\
	Do đó, tứ phân vị thứ nhất của mẫu số liệu ghép nhóm là \[Q'_1 =240+\dfrac{\dfrac{20}{4}-3}{7} (340-240)=\dfrac{1\,880}{7}.\] 
	Tứ phân vị thứ ba của mẫu  số liệu gốc là $\dfrac{1}{2} \left(x_{15}+x_{16} \right)\in [340;440)$. \\
	Do đó, tứ phân vị  thứ ba của mẫu số liệu ghép nhóm là \[Q_3 ^{{'}}=340+\dfrac{\dfrac{3.20}{4}-(3+7)}{7} (440-340)=\dfrac{2\,880}{7}.\]
	 \itemch Khoảng tứ phân vị của mẫu số liệu ghép nhóm là $\Delta_Q'=Q'_3-Q'_1=\dfrac{1\,000}{7}$.\\
	Khoảng biến thiên của mẫu số liệu ghép nhóm lớn  hơn mẫu số liệu; khoảng tứ phân vị của mẫu số liệu ghép nhóm nhỏ hơn mẫu số liệu.
	\end{itemchoice}
} 
\end{ex}

% nguồn: Lop_12_Dai_so_Chuong_3
\begin{ex}%[2D3H1-3]%[ma:63a47e85]
	Thống kê điểm trung bình môn Toán của một  số học sinh lớp $11$ được cho ở bảng sau
	\begin{center}
		\begin{tabular}{|c|c|c|c|c|c|c|c|}
			 \hline
			Khoảng điểm & $[6{,}5;7)$& $[7;7{,}5)$& $[7{,}5;8)$& $[8;8{,}5)$& $[8{,}5;9)$& $[9;9{,}5)$&  $[9{,}5;10)$\\
			 \hline
			Tần  số & $8$& $10$& $16$& $24$& $13$& $7$& $4$\\
			 \hline
		\end{tabular}
	\end{center}
	 \choiceTF
		{\True  Khoảng biến thiên của mẫu số liệu là $R=4$}
		{\True  Số trung bình của mẫu số liệu xấp xỉ bằng $8{,}12$}
		{Mốt  của mẫu số liệu là $6{,}21$}
		{Khoảng tứ phân vị của mẫu số  liệu ghép nhóm là $\Delta_Q=2{,}05$}
	 \loigiai{
		\begin{center}
			\begin{tabular}{|c|c|c|c|c|c|c|c|}
				\hline 
				Khoảng  điểm & $[6{,}5;7)$& $[7;7{,}5)$& $[7{,}5;8)$& $[8;8{,}5)$& $[8{,}5;9)$& $[9;9{,}5)$& $[9{,}5;10)$\\
				\hline 
				Giá trị đại diện & $6{,}75$& $7{,}25$& $7{,}75$& $8{,}25$& $8{,}75$&  $9{,}25$& $9{,}75$\\
				\hline 
				Tần số & $8$& $10$& $16$& $24$&  $13$& $7$& $4$\\
				\hline 
			\end{tabular}
		\end{center}
		\begin{itemchoice}
			 \itemch Khoảng biến thiên của mẫu số liệu là $R=4$.
			\itemch Số trung bình của mẫu số liệu xấp xỉ  bằng
			\[\left(6{,}75\cdot 8+7{,}25\cdot  10+7{,}75\cdot 16+8{,}25\cdot 24+8{,}75\cdot 13+9{,}25\cdot 7+9{,}75\cdot 4\right):82=8{,}12.
			\] 
			\itemch Nhóm chứa mốt của mẫu số liệu là $[8;8{,}5)$.\\ 
			Mốt  của mẫu số liệu là $M_0=8+\dfrac{24-16}{(24-16)+(24-13)} \cdot (8{,}5-8)=8{,}21$.
			\itemch Gọi $x_1;x_2;x_3;\ldots;x_{85}$ lần lượt là  tần số theo thứ tự không giảm.\\
			 Ta có
			\begin{itemize}
				\item $x_1, \ldots,  x_8 \in [6{,}5;7)$;
				\item  $x_9, \ldots, x_{18} \in [7;7{,}5)$;
				\item $x_{19}, \ldots, x_{34} \in  [7{,}5;8)$;
				\item $x_{35},  \ldots, x_{58} \in [8;8{,}5)$;
				\item $x_{59},  \ldots, x_{71} \in [8{,}5;9);\ldots$.
			\end{itemize}
			Tứ  phân vị thứ hai của dãy số liệu là $\dfrac{1}{2} \left(x_{41}+x_{42} \right)$ thuộc nhóm $[8;8{,}5)$ nên tứ phân vị thứ hai của mẫu số liệu là \[Q_2=8+\dfrac{\dfrac{82}{2}-34}{24} (8{,}5-8)=8{,}15.\]
			Tứ phân vị thứ nhất của dãy số liệu là $\dfrac{1}{2} \left(x_{20}+x_{21} \right)$ thuộc nhóm $[7{,}5;8)$ nên tứ phân vị thứ nhất của mẫu số liệu là \[Q_1=7{,}5+\dfrac{\dfrac{82}{4}-18}{16} (8-7{,}5)=7{,}58.\] 
			Tứ phân vị thứ ba của dãy số liệu là $\dfrac{1}{2} \left(x_{61}+x_{62} \right)$ thuộc nhóm $[8{,}5;9)$ nên tứ phân vị thứ ba của mẫu số liệu là \[Q_3=8{,}5+\dfrac{\dfrac{3\cdot 82}{4}-58}{13}  (9-8{,}5)=8{,}63.\]
			Khoảng tứ phân vị của mẫu số liệu ghép nhóm  là $\Delta_Q=1{,}05$.
		\end{itemchoice}
	} 
\end{ex}

% nguồn: Lop_12_Dai_so_Chuong_3
\begin{ex}%[2D3V1-3]%[ma:423bd7a0]
	Một công ty cung cấp nước sạch thống kê lượng nước các hộ gia đình trong một khu vực tiêu thụ trong một  tháng ở bảng sau
	\begin{center}
		\begin{tabular}{|c|c|c|c|c|c|}
			 \hline
			Lượng nước tiêu thụ $\left(\mathrm{m}^3\right)$  & $[3;6)$& $[6;9)$& $[9;12)$& $12;15)$& $15;18)$\\
			\hline 
			Số hộ gia đình & $24$& $57$& $42$& $29$& $8$\\ 
			\hline 
		\end{tabular}
	\end{center}
	\choiceTF 
		{\True Số trung bình của mẫu số liệu ghép nhóm  đã cho là $9{,}375$}
		{\True Mốt của mẫu số liệu  là $M_o=8{,}0625$}
		{\True Khoảng  biến thiên của mẫu số liệu ghép nhóm đã cho là $15$}
		{Công ty muốn gửi một thông báo khuyến nghị tiết kiệm nước đến $25\%$ các hộ gia đình có lượng nước tiêu thụ cao nhất. Khi đó công ty nên gửi thông báo tiết kiệm nước đến các hộ gia đình có lượng  nước tiêu thụ từ $14{,}79$\,m$^3$ nước trở lên}
	\loigiai{ 
		\begin{itemchoice}
			\itemch Cỡ mẫu $n=160$. Bảng  tần số ghép nhóm của mẫu số liệu trên như sau
			\begin{center}
				\begin{tabular}{|c|c|c|c|c|c|}
					 \hline
					Lượng  nước tiêu thụ $\left(\mathrm{m}^3\right)$ & $[3;6)$& $[6;9)$& $[9;12)$& $[12;15)$& $[15;18)$\\
					 \hline
					Giá trị đại diện & $4{,}5$& $7{,}5$& $10{,}5$& $13{,}5$&  $16{,}5$\\
					\hline 
					Số hộ gia  đình & $24$& $57$& $42$& $29$& $8$\\
					 \hline
				\end{tabular}
			\end{center}
			Số trung bình của mẫu số liệu  ghép nhóm đã cho là
			\[\overline{x}=\dfrac{1}{160} \left(24\cdot 4{,}5+57\cdot 7{,}5+42\cdot 10{,}5+29\cdot  13{,}5+8\cdot 16{,}5\right)=9{,}375.
			 \]
			\itemch Khoảng biến thiên của mẫu số liệu là  $R=18-3=15$.
			\itemch Nhóm chứa mốt của mẫu số liệu ghép  nhóm trên là nhóm $\left[6;9\right)$.\\
			 Do đó $u_m=6$; $n_m=57$; $n_{m-1}=24$; $n_{m+1}=42$; $u_{m+1}=9$.\\
			Mốt của mẫu số liệu  là \[M_o=6+\dfrac{(57-24)}{(57-24)+(57-42)} \cdot (9-6)=8{,}0625.\]
			\itemch Gọi $x_1;x_2;\ldots;x_{160}$ là mẫu số liệu  được xếp theo thứ tự không giảm.\\
			Ta  có
			\begin{itemize}
				\item $x_1, \ldots, x_{24} \in  [3;6)$;
				\item $x_{25}, \ldots, x_{81} \in  [6;9)$;
				\item $x_{82}, \ldots, x_{123} \in [9;12)$; 
				 \item $x_{124}, \ldots, x_{152} \in [12;15)$;
				 \item $x_{153}, \ldots, x_{160} \in [15;18)$.
			\end{itemize}
			Cỡ mẫu $n=160$ là số chẵn nên  trung vị là $M_e=\dfrac{1}{2} \left(x_{80}+x_{81} \right)$.\\
			Trung vị của mẫu số liệu  là $M_e=6+\dfrac{\dfrac{160}{2}-24}{57} \left(9-6\right)\approx 8{,}95$.\\
			$25\%$ các  hộ gia đình có lượng nước tiêu thụ cao nhất có lượng nước tiêu thụ không nhỏ hơn $Q_3$, với $Q_3$ là tứ phân vị thứ ba của mẫu số liệu.\\
			Tứ phân vị thứ ba của dãy số liệu $x_1;x_2;\ldots;x_{160}$ là  $\dfrac{1}{2} \left(x_{120}+x_{121} \right)$.\\
			Do $x_{120}$ và $x_{121}$ thuộc nhóm $[9;12)$  nên tứ phân vị thứ ba của mẫu số liệu là
			\[Q_3=9+\dfrac{\dfrac{160\cdot 3}{4}-(24+57)}{42} \cdot  (12-9)\approx 11{,}79.
			\] 
			Vậy công ty nên gửi thông báo tiết kiệm nước đến các hộ gia đình có lượng nước tiêu  thụ từ $11{,}79$\,m$^3$ nước trở lên.
		\end{itemchoice}
	 }
\end{ex}

% nguồn: Lop_12_Dai_so_Chuong_3
\begin{ex}%[2D3V1-3]%[ma:b7dbf60e]
	Hằng ngày ông Thắng đều đi xe buýt từ nhà đến cơ quan. Dưới đây là bảng thống kê thời gian của $100$ lần ông  Thắng đi xe buýt từ nhà đến cơ quan.
	\begin{center}
		\begin{tabular}{|c|c|c|c|c|c|c|}
			 \hline
			Thời gian  (phút)& $[15;18)$& $[18;21)$& $[21;24)$& $[24;27)$& $[27;30)$& $[30;33)$\\
			 \hline
			Số  lượt& $22$& $38$& $27$& $8$& $4$& $1$\\
			 \hline
		\end{tabular}
	\end{center}
	\choiceTF 
		{\True Cỡ mẫu  $n=100$}
		{Tứ phân vị thứ nhất của mẫu số liệu ghép nhóm  là $Q_1=\dfrac{683}{38}$}
		{Khoảng tứ phân vị của mẫu số liệu ghép là  $\Delta_Q=\dfrac{515}{114}$}
		{\True Biết trong $100$ lần đi trên, chỉ có  đúng một lần ông Thắng đi hết hơn $29$ phút. Thời gian của lần đó là giá trị ngoại lệ của mẫu số liệu ghép nhóm}
	 \loigiai{
		\begin{itemchoice}
			\itemch Cỡ mẫu  $n=100$.\\
			Gọi $x_1;x_2;\ldots;x_{100}$ là  mẫu số liệu gốc gồm thời gian 100 lần đi xe buýt của ông Thắng.\\
			 Ta có
			\begin{itemize}
				\item $x_1, \ldots, x_{22}  \in [15;18)$;
				\item $x_{23}, \ldots, x_{60} \in  [18;21)$;
				\item $x_{61}, \ldots, x_{87} \in  [21;24)$;
				\item $x_{88}, \ldots, x_{95} \in  [24;27)$;
				\item $x_{96}, \ldots,  x_{99} \in [27;30)$;
				\item $x_{100} \in  [30;33)$.
			\end{itemize}
			\itemch Tứ  phân vị thứ nhất của mẫu số liệu gốc là $\dfrac{1}{2} \left(x_{25}+x_{26} \right)\in [18;21)$. \\
			Do đó, tứ phân vị thứ nhất của mẫu  số liệu ghép nhóm là \[Q_1=18+\dfrac{\dfrac{100}{4}-22}{38} \left(21-18\right)=\dfrac{693}{38}.\]
			\itemch Tứ phân vị thứ ba của  mẫu số liệu gốc là $\dfrac{1}{2} \left(x_{75}+x_{76} \right)\in [21;24)$. \\
			Do đó, tứ phân vị thứ ba  của mẫu số liệu ghép nhóm là \[Q_3=21+\dfrac{\dfrac{3.100}{4}-\left(22+38\right)}{27} \left(24-21\right)=\dfrac{68}{3}.\]
			Vậy khoảng tứ phân vị của mẫu  số liệu ghép nhóm là $\Delta_Q=\dfrac{68}{3}-\dfrac{693}{38}=\dfrac{505}{114}$
			\itemch Trong lần duy nhất ông Thắng đi  hết hơn $29$ phút, thời gian đi của ông thuộc nhóm $[30;33)$. \\
			Vì $Q_3+1,5\Delta_Q=\dfrac{6683}{228} < 30$ nên thời gian của lần ông Thắng đi  hết hơn $29$ phút là giá trị ngoại lệ của mẫu số liệu ghép nhóm.
		\end{itemchoice}
	 }
\end{ex}

% nguồn: Lop_12_Dai_so_Chuong_3
\begin{ex}%[2D3V1-3]%[ma:f3e0c866]
	Giả sử kết quả khảo sát hai khu vực A và B về độ  tuổi kết hôn của một số phụ nữ vừa lập gia đình được cho ở bảng sau
	\begin{center}
		\begin{tabular}{|c|c|c|c|c|c|}
			\hline 
			 Tuổi kết hôn & $[19;22)$& $[22;25)$& $[25;28)$& $[28;31)$& $[31;34)$\\
			 \hline
			Số phụ nữ khu vực A& $10$& $72$& $31$& $25$& $7$\\ 
			 \hline
			Số phụ  nữ khu vực B& $47$& $40$& $11$& $2$& $0$\\
			\hline 
		\end{tabular}
	\end{center}
	 \choiceTF
		{\True Khoảng biến thiên của mẫu số liệu ghép nhóm ứng với khu  vực A là $15$ (tuổi)}
		{\True Khoảng biến thiên của mẫu số liệu ghép nhóm ứng với khu  vực B là $12$ (tuổi)}
		{Khoảng tứ phân vị của mẫu số liệu  ghép ứng với khu vực A là $\Delta_Q=\dfrac{61}{3}$ (tuổi)}
		{\True Nếu so sánh theo khoảng tứ phân vị thì phụ nữ ở khu vực B có độ  tuổi kết hôn đồng đều hơn}
	 \loigiai{
		\begin{itemchoice}
			\itemch Khoảng biến thiên của mẫu số liệu  ghép nhóm ứng với khu vực A là $34-19=15$ (tuổi).
			\itemch Khoảng biến thiên của mẫu số liệu ghép nhóm ứng với khu vực B là $31-19=12$  (tuổi).
			\itemch Cỡ mẫu $n=100$.\\ 
			Gọi $x_1;x_2;\ldots;x_{100}$ là mẫu số liệu gốc về độ tuổi kết hôn của phụ nữ ở khu vực A  được xếp theo thứ tự không giảm.\\
			Ta  có
			\begin{itemize}
				 \item $x_1;x_2;\ldots;x_{10} \in [19;22)$;
				\item  $x_{11};\ldots;x_{37} \in [22;25)$;
				 \item $x_{38};\ldots;x_{68} \in [25;28)$;
				\item $x_{69};\ldots;x_{93}  \in [28;31)$;
				 \item $x_{94};\ldots;x_{100} \in [31;34)$.
			\end{itemize}
			Tứ phân vị thứ nhất của mẫu số liệu  gốc là $\dfrac{1}{2} \left(x_{25}+x_{26} \right)\in [22;25)$. \\
			Do đó, tứ phân  vị thứ nhất của mẫu số liệu ghép nhóm là \[Q_1=22+\dfrac{\dfrac{100}{4}-10}{27} (25-22)=\dfrac{71}{3}.\]
			Tứ phân vị thứ  ba của mẫu số liệu gốc là $\dfrac{1}{2} \left(x_{75}+x_{76} \right)\in [28;31)$. \\
			Do đó, tứ phân vị thứ ba của mẫu số  liệu ghép nhóm là \[Q_3=28+\dfrac{\dfrac{3\cdot 100}{4}-(10+27+31)}{25} (31-28)=\dfrac{721}{25}.\]
			Khoảng tứ phân vị của mẫu số  liệu ghép nhóm là $\Delta_Q=Q_3-Q_1=\dfrac{388}{75}$.\\
			Gọi $y_1;y_2;\ldots;y_{100}$ là mẫu số liệu gốc về độ tuổi kết hôn của phụ nữ ở khu vực B được xếp theo thứ  tự không giảm.\\
			Ta  có
			\begin{itemize}
				\item  $y_1;y_2;\ldots;y_{47} \in [19;22)$;
				\item $y_{48};\ldots;y_{87} \in [22;25)$; 
				\item  $y_{88};\ldots;y_{98} \in [25;30)$;
				\item $y_{99};y_{100} \in [28;31)$. 
			\end{itemize}
			Tứ phân vị thứ  nhất của mẫu số liệu gốc là $\dfrac{1}{2} \left(y_{25}+y_{26} \right)\in [19;22)$. \\
			Do đó, tứ phân vị thứ nhất của mẫu số liệu ghép nhóm  là \[Q'_1 =19+\dfrac{\dfrac{100}{4}}{47} (22-19)=\dfrac{968}{47}.\]
			Tứ phân vị thứ ba của mẫu số liệu gốc là $\dfrac{1}{2} \left(y_{75}+y_{76} \right)\in  [22;25)$. \\
			Do đó, tứ phân vị thứ ba của mẫu số liệu ghép  nhóm là \[Q'_3 =22+\dfrac{\dfrac{3\cdot 100}{4}-47}{40} (25-22)=\dfrac{241}{10}.\]
			\itemch Có $\Delta_Q' < \Delta_Q$ nên phụ  nữ ở khu vực B có độ tuổi kết hôn đồng đều hơn.
		\end{itemchoice}
	 }
\end{ex}

% nguồn: Lop_12_Dai_so_Chuong_3
\begin{ex}%[2D3V1-3]%[ma:771d206b]
	Tìm hiểu thời gian sử dụng điện thoại trong tuần đầu tháng  $6/2024$ của kỳ nghỉ hè lớp chủ nhiệm. Giaó viên chủ nhiệm thu được kết quả sau
	\begin{center}
		\begin{tabular}{|c|c|c|c|c|c|c|}
			 \hline
			Thời gian (giờ) & {$[0; 5)$} & {$[5; 10)$} & {$[10; 15)$} & {$[15; 20)$} & {$[20; 25)$} &  {$[25; 30]$} \\
			 \hline
			Số học  sinh & 2 & 6 & 8 & 9 & 3 & 2 \\
			\hline 
		\end{tabular}
	\end{center}
	Khi đó 
	\choiceTF 
		{Khoảng biến thiên của  mẫu số liệu ghép nhóm này là $25$}
		{\True Nhóm chứa  tứ phân vị thứ 3 là $[15; 20)$}
		{\True Số trung bình của thống  kê là $10$}
		{Khoảng tứ phân của mẫu số  liệu ghép nhóm này lớn hơn $10$}
	\loigiai{ 
		\begin{center}
			\begin{tabular}{|c|c|c|c|c|c|c|}
				\hline 
				Thời gian (giờ) & {$[0; 5)$} & {$[5; 10)$} & {$[10;  15)$} & {$[15; 20)$} & {$[20; 25)$} & {$[25; 30]$} \\
				 \hline
				Số học sinh & $2$ & $6$ & $8$  & $9 $& $3$ & $2$ \\
				 \hline
				Giá trị  đại diện & $2{,}5$ & $7{,}5$ & $12{,}5$ & $17{,}5$ & $22{,}5$ & $27{,}5$ \\
				\hline 
			\end{tabular}
		\end{center}
		\begin{itemchoice}
			\itemch Sai. Khoảng  biến thiên của mẫu số liệu ghép nhóm là $R=30-0=30$.
			\itemch Đúng. Vì $16< \dfrac{3n}{4}=\dfrac{3\cdot30}{4}=\dfrac{90}{4}=22{,}5< 25$ nên nhóm chứa  tứ phân vị thứ $3$ là $[15; 20)$.
			\itemch  Đúng. Thời gian sử dụng điện thoại trung bình của học sinh là
			$$\bar{x}=\dfrac{2\cdot2{,}5+6\cdot7{,}5+8\cdot12{,}5+9\cdot17{,}5+3\cdot22{,}5+2\cdot27{,}5}{30}=\dfrac{43}{3} \approx 14{,}3.$$ 
			\itemch Sai. Ta có $\dfrac{n}{4}=7{,}5$; $\dfrac{n}{2}=15$;  $\dfrac{3n}{4}=22{,}5$.\\
			$Q_1=5+\dfrac{\dfrac{30}{4}-2}{6}\cdot5=9{,}58$;  $Q_3=15+\dfrac{\dfrac{90}{4}-16}{9}\cdot5=\dfrac{335}{18} \approx 18{,}61$\\
			 $\Rightarrow \Delta_Q=Q_3-Q_1=\dfrac{325}{36} \approx 9{,}03< 10$.
		\end{itemchoice}
	} 
\end{ex}

% nguồn: Lop_12_Dai_so_Chuong_3
\begin{ex}%[2D3V1-3]%[ma:730aaaf8]
	Kết quả  đo chiều cao của $100$ cây keo $3$ năm tuổi tại một nông trường được cho ở bảng sau
	\begin{center}
		\begin{tabular}{|c|c|c|c|c|c|}
			 \hline
			Chiều cao (m)  & $[8{,}4;8{,}6)$& $[8{,}6;8{,}8)$& $[8{,}8;9{,}0)$& $[9{,}2;9{,}4)$& $[9{,}0;9{,}2)$\\
			 \hline
			Số cây & $5$& $12$&  $25$& $44$& $14$\\
			\hline 
		\end{tabular}
	\end{center}
	\choiceTF 
		{\True Khoảng  biên thiên của mẫu số liệu ghép nhóm $1$\,m}
		{\True Tứ  phân vị thứ ba của mẫu số liệu ghép nhóm là $Q_3=10{,}5$}
		{Khoảng tứ phân vị của mẫu số liệu ghép nhóm  là $\Delta_Q=2{,}06$}
		{Trong $100$ cây keo trên có $1$  cây cao $8{,}4$\,m thì chiều cao của cây keo này là giá trị ngoại lệ}
	\loigiai{ 
		\begin{itemchoice}
			\itemch Khoảng biên thiên của mẫu số  liệu ghép nhóm $9{,}4-8{,}4=1$\,(m).
			 \itemch Cỡ mẫu $n=100$;\\
			 Gọi $x_1;x_2;\ldots;x_{100}$ là mẫu số liệu gốc được xếp theo thứ tự không giảm.\\
			Ta có 
			\begin{itemize}
				 \item $x_1;\ldots;x_5 \in [8{,}4;8{,}6)$;
				\item  $x_6;\ldots;x_{17} \in [8{,}6;8{,}8)$;
				\item $x_{18};\ldots;x_{42}  \in [8{,}8;9{,}0)$;
				\item $x_{43};\ldots;x_{86} \in  [9{,}0;9{,}2)$;
				\item $x_{87};\ldots;x_{100} \in [9{,}2;9{,}4)$. 
			\end{itemize}
			Tứ  phân vị thứ nhất của mẫu số liệu gốc là $\dfrac{1}{2} \left(x_{25}+x_{26} \right)\in [8{,}8;9{,}0)$.\\
			Do đó, tứ phân vị thứ nhất của mẫu số liệu ghép nhóm là \[Q_1=8{,}8+\dfrac{\dfrac{100}{4}-(5+12)}{25}  (9{,}0-8{,}8)=9{,}44.\]
			Tứ phân vị thứ ba của mẫu số liệu gốc là $\dfrac{1}{2} \left(x_{75}+x_{76} \right)\in [9{,}0;9{,}2)$.\\ 
			Do đó, tứ phân vị thứ ba của mẫu số  liệu ghép nhóm là \[Q_3=9{,}0+\dfrac{\dfrac{3\cdot 100}{4}-\left(5+12+25\right)}{44} \left(9{,}2-9{,}0\right)=10{,}5.\]
			\itemch Khoảng tứ phân vị của mẫu số liệu ghép nhóm  là $\Delta_Q=Q_3-Q_1=1{,}06$.
			\itemch Giá trị $x$ trong mẫu số liệu là giá  trị ngoại lệ nếu $x > Q_3+1{,}5\Delta_Q$ hoặc $x < Q_1-1{,}5\Delta_Q$.\\
			Hay $x > 10{,}5+1{,}5\cdot 1{,}06=12{,}09$ hoặc $x < 9{,}44-1{,}5\cdot 1{,}06=7{,}85$.\\ 
			 Vậy cây cao $8{,}4$\,m không phải là giá trị ngoại lệ.
		\end{itemchoice}
	 }
\end{ex}


%%%%%%%%%%%%%%%% TRẢ LỜI NGẮN %%%%%%%%%%%%%%%%

% nguồn: Lop_12_Dai_so_Chuong_3
\begin{bt}%[2D3H1-3]%[ma:41c3a775]
	Thời gian sử dụng điện thoại trong một  ngày của 30 sinh viên được ghi lại ở bảng sau (đơn vị: phút).
	\begin{center}
		\begin{tabular}{|c|c|c|c|c|c|}
			 \hline
			Thời gian (phút) & $[0;60)$& $[60;120)$& $[120;180)$&  $[180;240)$& $[240;300)$\\
			\hline 
			Số sinh viên & $2$& $7$&  $7$& $10$& $4$\\
			 \hline
		\end{tabular}
	\end{center}
	Hãy tìm khoảng tứ  phân vị của mẫu số liệu ghép nhóm trên.
	 \shortans{112}
	\loigiai{ 
		Gọi $x_1;x_2;\ldots;x_{30}$  là mẫu số liệu được xếp theo thứ tự không giảm.\\
		Ta có 
		\begin{itemize}
			\item $x_1, x_2 \in  [0;60)$;
			\item $x_3, \ldots,  x_9 \in [60;120)$;
			\item $x_{10}, \ldots, x_{16} \in  [120;180)$;
			\item  $x_{17}, \ldots, x_{26} \in [180;240)$;
			 \item $x_{27}, \ldots, x_{30} \in [240;300)$.
		\end{itemize}
		Tứ phân vị thứ  hai của mẫu số liệu ghép nhóm là $\dfrac{1}{2} \left(x_{15}+x_{16} \right)$.\\
		Do $x_{15}$ và $x_{16}$ thuộc nhóm $[120;180)$ nên tứ phân vị thứ hai của mẫu số liệu ghép  nhóm là
		\[Q_2=120+\dfrac{\dfrac{30}{2}-\left(2+7\right)}{7} \left(180-120\right)=\dfrac{1\,200}{7} \approx 171{,}43.\] 
		Tứ phân vị thứ nhất của mẫu  là $x_8 \in [60;120)$. \\
		Do đó, tứ phân vị thứ  nhất của mẫu số liệu ghép nhóm là \[Q_1=60+\dfrac{\dfrac{30}{4}-2}{7} \left(120-60\right)=\dfrac{750}{7} \approx 107{,}14.\]
		Tứ phân vị thứ ba của mẫu số  liệu là $x_{23} \in [180;240)$.\\
		Do đó, tứ phân vị  thứ ba của mẫu số liệu ghép nhóm là
		\[Q_3=180+\dfrac{\dfrac{3\cdot 30}{4}-\left(2+7+7\right)}{10}  \left(240-180\right)=219.\]
		Khoảng tứ phân vị của mẫu  số liệu ghép nhóm là $\Delta_Q=111{,}86$.
	 }
\end{bt}

% nguồn: Lop_12_Dai_so_Chuong_3
\begin{bt}%[2D3H1-3]%[ma:3f4169bd]
	Cho  mẫu số liệu ghép nhóm về thời gian (phút) đi từ nhà đến nơi làm việc của các nhân viên của một công ty như sau
	\begin{longtable}{|c|c|c|c|c|c|c|c|}
		\hline 
		Thời gian &$[15;20)$&$[20;25)$ &$[25;30)$&$[30;35)$&  $[35;40)$&$[40;45)$& $[45;50)$ \\
		\hline 
		Số nhân viên &$6$& $14$&$30$&$25$&$22$&  $15$&$8$ \\
		\hline 
	\end{longtable}
	Tìm khoảng tứ phân  vị của mẫu số liệu trên.
	\shortans{11,7} 
	\loigiai{Cỡ  mẫu là $n=6+14+30+25+22+15+8=120$.\\
		Gọi $x_1$, $x_2$,$\ldots$, $x_{120}$ là thời gian đi từ nhà đến nơi làm việc của các nhân viên xếp theo thứ tự không  giảm.\\
		Ta có $x_1$, $x_2$,$\ldots$, $x_{6} \in [15;20)$, $x_{7}$, $x_{8}$,$\ldots$, $x_{20} \in [20;25)$, $x_{21}$, $x_{22}$,$\ldots$, $x_{50} \in [25;30)$, $x_{51}$, $x_{52}$,$\ldots$, $x_{75} \in  [30;35)$, $x_{76}$, $x_{77}$,$\ldots$, $x_{97} \in [35;40)$, $x_{98}$, $x_{98}$,$\ldots$, $x_{112} \in [40;45)$, $x_{113}$, $x_{114}$,$\ldots$, $x_{120} \in [45;50)$.
		\\ 
		Tứ phân vị thứ hai $Q_2$ của mẫu số liệu là  giá trị $\dfrac{x_{60}+x_{61}}{2}$.\\
		Tứ phân vị thứ nhất $Q_1$ của mẫu số liệu là giá trị $\dfrac{x_{30}+x_{31}}{2}$.  Nhóm chứa tứ phân vị thứ nhất là $[25;30)$.\\
		Vậy $Q_1=25+\dfrac{\dfrac{120}{4}-(6+14)}{30}  \cdot \left(30-25\right)=\dfrac{80}{3}$.\\
		 Tứ phân vị thứ ba $Q_3$ của mẫu số liệu là $\dfrac{x_{79}+x_{80}}{2}$. Nhóm chứa tứ phân vị thứ ba là $[35;40)$.\\
		Vậy $Q_3=35+\dfrac{\dfrac{3  \cdot 120}{4}-(6+14+30+25)}{22} \cdot (40-35)=\dfrac{845}{22}$.\\
	Vậy khoảng tứ phân vị là $Q_3-Q_1=\dfrac{845}{22}-\dfrac{80}{3}=\dfrac{775}{66}\approx 11{,}7$.} 
\end{bt}

% nguồn: Lop_12_Dai_so_Chuong_3
\begin{bt}%[2D3H1-3]%[ma:96b3ac27]
	Bảng dưới biểu diễn mẫu số liệu ghép nhóm vể chiều cao của $42$ mẫu cây  ở một vườn thực vật (đơn vị: centimét).
	\begin{center}
		\begin{tabular}{|c|c|c|c|c|c|c|}
			\hline 
			Nhóm & $[40;45)$& $[45;50)$& $[50;55)$& $[55;60)$&  $[60;65)$& $[65;70)$\\
			\hline 
			Tần số & $5$& $10$&  $7$& $9$& $7$& $4$\\
			\hline 
		\end{tabular}
	\end{center}
	Tính khoảng tứ phân vị của mẫu số liệu ghép nhóm đó (làm tròn kết quả đến hàng phẩn  mười).
	\shortans{12,7} 
	\loigiai{ 
		Số phần  tử của mẫu là $n=42$.\\
		Ta  có $\dfrac{n}{4}=\dfrac{42}{4}=10{,}5$. \\
		Mà $5< 10{,}5< 15$ nên nhóm $2$ là nhóm đầu tiên có tần số  tích luỹ lớn hơn hoặc bằng $10{,}5$. \\
		Xét nhóm $2$ là nhóm $[45;50$) có $s=45$; $h=5$; $n_2=10$ và nhóm $1$ là nhóm  $[40;45)$ có $cf_1=5$.\\
		Áp dụng công thức, ta có tứ phân vị thứ nhất là $Q_1=45+\left(\dfrac{10{,}5-5}{10} \right)\cdot  5=47{,}75$\,(cm).\\
		Ta có $\dfrac{3n}{4}=\dfrac{3\cdot  42}{4}=31{,}5$.\\
		Mà $31< 31{,}5< 38$ nên nhóm $5$ là nhóm đầu  tiên có tần số tích luỹ lớn hơn hoặc bằng $31{,}5$.\\
		Xét nhóm $5$ là nhóm  $[60;65)$ có $t=60$; $l=5$; $n_5=7$ và nhóm $4$ là nhóm $[55;60)$ có $cf_4=31$.\\
		Áp dụng công thức, ta có tứ phân vị  thứ ba là $Q_3=60+\left(\dfrac{31{,}5-31}{7} \right)\cdot 5\approx 60{,}4$\,( cm).\\
		Vậy khoảng tứ phân vị của mẫu số liệu ghép nhóm  đã cho là \[\Delta_Q=Q_3-Q_1 \approx 60{,}4-47{,}75=12{,}65\,\mathrm{(cm)}.\]
	} 
\end{bt}

% nguồn: Lop_12_Dai_so_Chuong_3
\begin{bt}%[2D3H1-3]%[ma:2873fdd1]
	 Bảng sau thống kê cân nặng của $50$ quả xoài Thanh Ca được lựa chọn ngẫu nhiên sau khi thu hoạch ở một nông trường.
	\begin{center}
		\begin{tabular}{|c|c|c|c|c|c|}
			 \hline
			Cân nặng (gam) & $[250;290)$& $[290;330)$& $[330;370)$& $[370;410)$& $[410;450)$\\ 
			\hline 
			Số quả xoài & $3$& $13$& $18$& $11$& $5$\\ 
			\hline 
		\end{tabular}
	\end{center}
	Khoảng tứ phân  vị của mẫu số liệu ghép nhóm đó bằng $\dfrac{a}{b}$ với $\dfrac{a}{b}$ là phân số tối giản và $a,b\in \mathbb{Z}$. Tính giá trị biểu thức $T=a+b$.
	\shortans{9223} 
	\loigiai{ 
		Cỡ mẫu $n=50$.\\ 
		Gọi $x_1;x_2;\ldots;x_{50}$ là mẫu số liệu gốc gồm cân  nặng của $50$ quả xoài được xếp theo thứ tự không giảm.\\
		 Ta có
		\begin{itemize}
			 \item $x_1, x_2, x_3 \in [250;290)$;
			\item  $x_4, \ldots, x_{16} \in [290;330)$;
			\item $x_{17}, \ldots, x_{34}  \in [330;370)$;
			\item $x_{35}, \ldots, x_{45} \in [370;410)$; 
			\item $x_{46}, \ldots,  x_{50} \in [410;450)$.
		\end{itemize}
		Tứ  phân vị thứ nhất của mẫu số liệu gốc là $x_{13} \in [290;330)$. \\
		Do  đó, tứ phân vị thứ nhất của mẫu số liệu ghép nhóm là
		\[Q_1=290+\dfrac{\dfrac{50}{4}-3}{13}  \left(330-290\right)=\dfrac{4150}{13}.\]
		Tứ phân vị thứ ba  của mẫu số liệu gốc là $x_{38} \in [370;410)$. \\
		Do đó, tứ phân vị thứ ba  của mẫu số liệu ghép nhóm là
		\[Q_3=370+\dfrac{\dfrac{3\cdot  50}{4}-\left(3+13+18\right)}{11} \left(410-370\right)=\dfrac{4\,210}{11}.\]
		Vậy khoảng tứ phân vị của mẫu số liệu ghép nhóm là  $\Delta_Q=\dfrac{4210}{11}-\dfrac{4\,150}{13}=\dfrac{9\,080}{143}$.\\
		Vậy giá  trị biểu thức $T=a+b=9\,080+143=9\,223$.
	} 
\end{bt}

% nguồn: Lop_12_Dai_so_Chuong_3
\begin{bt}%[2D3H1-3]%[ma:34a8177a]
	Kết quả điều tra về số  giờ làm thêm trong một tuần của $100$ sinh viên được cho cho như bảng sau
	\begin{center}
		\begin{tabular}{|c|c|c|c|c|c|}
			\hline 
			 Số giờ làm thêm & $[2;4)$& $[4;6)$& $[6;8)$& $[8;10)$& $[10;12)$\\
			 \hline
			Số sinh  viên & $12$& $20$& $37$& $21$& $10$\\
			 \hline
		\end{tabular}
	\end{center}
	Tìm khoảng tứ phân  vị của số liệu đó.
	\shortans{3,27} 
	 \loigiai{
		Cỡ mẫu  $n=100$.\\
		Gọi $x_1;x_2;\ldots;x_{100}$ là mẫu số liệu được xếp theo  thứ tự không giảm.\\
		Tứ phân vị thứ hai của dãy số liệu $x_1;x_2;\ldots;x_{100}$ là $\dfrac{1}{2}  \left(x_{50}+x_{51} \right)$.\\
		Do $x_{50}$ và $x_{51}$ thuộc nhóm $[6;8)$ nên \[Q_2=6+\dfrac{\dfrac{2.100}{4}-32}{37}\cdot\left(8-6\right)\approx  6{,}97.\]
		Tứ phân vị thứ nhất của dãy số liệu $x_1;x_2;\ldots;x_{100}$ là  $\dfrac{1}{2} \left(x_{25}+x_{26} \right)$.\\
		Do $x_{25}$ và $x_{26}$  thuộc nhóm $[4;6)$ nên $Q_1=4+\dfrac{\dfrac{1.100}{4}-12}{20}.\left(6-4\right)=5{,}3$.\\
		Tứ phân vị thứ  ba của dãy số liệu $x_1;x_2;\ldots;x_{100}$ là $\dfrac{1}{2} \left(x_{75}+x_{76} \right)$.\\
		Do $x_{75}$  và $x_{76}$ thuộc nhóm $[8;10)$ nên \[Q_3=8+\dfrac{\dfrac{3\cdot 100}{4}-69}{21}\cdot\left(10-8\right)\approx 8{,}57.\]
		Khoảng tứ phân vị là $\Delta Q=3{,}27$. 
	 }
\end{bt}

% nguồn: Lop_12_Dai_so_Chuong_3
\begin{bt}%[2D3H1-3]%[ma:aad3a28a]
	Biểu đồ dưới đây biểu diễn số lượt khách hàng đặt  bàn qua hình thức trực tuyến mỗi ngày trong quý III năm $2\,022$ của một nhà hàng. Cột thứ nhất biểu diễn số ngày có từ $1$ đến dưới $6$ lượt đặt bàn; cột thứ hai biểu diễn số ngày có từ 6 đến dưới $11$ lượt đặt bàn; $\ldots$.
	\begin{center}
		\begin{tikzpicture}[
					join = round,cap = round, 
					declare function={ 
					 x=2;
					y=1; 
					kc=1; 
					 },
					 xscale = 1/x,yscale = 1/y,
					font = \footnotesize] 
				\foreach \x/\y[count =  \i from 1] in {4/1.4,5/3,6/2.5,7/1.8,8/0.5}{
					 \pgfmathsetmacro{\j}{kc*\i}
					 \pgfmathsetmacro{\realy}{10*\y}
					\fill[pattern=north  east lines,draw,black] (\i+1,0) rectangle +(1,\y);
					\path (\x-1.5,\y+0.25) node{$\pgfmathprintnumber{\realy}$}; 
					\global\let\n=\j 
				 }
				\draw[<->] (\n+1.75*x,0) -| (0,3.25+y); 
				\foreach \y  in {0,.5,1,1.5,2,2.5,3,3.5}{
					 \pgfmathsetmacro{\realy}{10*\y}
					\draw (.1*x,\y)--(-.1*x,\y)  node[left]{$\pgfmathprintnumber{\realy}$};
				 }
				 \foreach \x in{1,6,11,16,21,26}{
					\path  (.2*\x+2,0) node[below]{\x};
				 }
				\path (6.5,4+y) node[text width=10cm,align=center]{\bf{Số lượt khách đặt  bàn trực tuyến}};
				 \path (\n+1.75*x+3,0) node{Số lượt đặt bàn};
				\path (0,3.5+y) node[left]{Số  ngày};
			\end{tikzpicture}
	\end{center}
	Hãy tìm  khoảng tứ phân vị của mẫu số liệu ghép nhóm cho bởi biểu đồ trên.
	\shortans{8,5} 
	\loigiai{ 
		Cỡ  mẫu $n=92$;
		Gọi $x_1;x_2;\ldots;x_{92}$ là mẫu số liệu gốc về số  lượt khách hàng đặt bàn qua hình thức trực tuyến mỗi ngày trong quý III năm $2\,022$ của nhà hàng được xếp theo thứ tự không giảm.\\
		Ta có 
		\begin{itemize}
			\item $x_1;\ldots;x_{14} \in [1;6)$; 
			\item $x_{15};\ldots;x_{44} \in  [6;11)$;
			\item  $x_{45};\ldots;x_{69} \in [11;16)$;
			\item  $x_{70};\ldots;x_{87} \in [16;21)$;
			 \item $x_{88};\ldots;x_{92} \in [21;26)$.
		\end{itemize}
		Tứ phân vị thứ nhất của mẫu số  liệu gốc là $\dfrac{1}{2} \left(x_{23}+x_{24} \right)\in [6;11)$.\\ Do đó, tứ phân vị thứ nhất của mẫu số liệu ghép nhóm là \[Q_1=6+\dfrac{\dfrac{92}{4}-14}{30} \left(11-6\right)=7{,}5.\]
		 Tứ phân vị thứ ba của mẫu số liệu gốc là $\dfrac{1}{2} \left(x_{69}+x_{70} \right)\in [11;16)$ và $[16;21)$.\\
		Do đó, tứ phân vị thứ ba của mẫu số  liệu ghép nhóm là $Q_3=16$.\\
		Khoảng tứ phân vị của mẫu số liệu ghép nhóm là  $\Delta_Q=Q_3-Q_1=8{,}5$.
	} 
\end{bt}

% nguồn: Lop_12_Dai_so_Chuong_3
\begin{bt}%[2D3H1-3]%[ma:1560c5d1]
	\immini { 
	Biểu đồ biểu diễn số lượt khách hàng đặt bàn qua hình thức trực tuyến mỗi ngày trong quý III năm 2022 của một nhà hàng. Cột thứ nhất biểu diễn số ngày có từ $1$ đến dưới $6$ lượt đặt bàn; cột thứ hai biểu diễn số ngày có từ $6$ đến dưới $11$  lượt đặt bàn; \ldots\\
	Hãy tìm  khoảng tứ phân vị của mẫu số liệu ghép nhóm cho bởi biểu đồ trên.
	 }{
	\begin{tikzpicture}line join=round,line cap=round,>=stealth,x=0.8cm,y=0.4cm]
			 \tikzset{ethicknode/.style={scale=0.8}}
			\draw[->] (0,0)--(0,8) node[left]{Số  ngày};
			\draw[->]  (0,0)--(7,0) node[above,align=left]{Số lượt\\ đặt bàn};
			\foreach \y in {0,5,10,...,35}{ 
				\draw[dotted,gray!60]  (0,\y/5)--++(6,0);
				\path  (0,\y/5) node[left]{$\y$};
				 };
			\foreach \x  in {1,6,...,26}{
				\draw  ({1+(\x-1)/5},0) node[below]{$\x$};}
			\foreach \a/\b  in {1/14,2/30,3/25,4/18,5/5}{
				\draw[pattern=crosshatch dots] (\a,0)  rectangle (\a+1,{\b/5}) node[shift={(150:0.8)}]{$\b$};
			} 
			\path (current bounding box.north) node{\bf Số lượt khách hàng đặt  bàn trực tuyến};
		\end{tikzpicture}
	} 
	\shortans{8,5} 
	 \loigiai{
		Cỡ mẫu $n=92$;\\ 
		$Q_1=6+\dfrac{\dfrac{92}{4}-14}{30} (11-6)=7{,}5$; $Q_3=11+\dfrac{\dfrac{3\cdot 92}{4}-44}{25}  (16-11)=16$.\\
		Khoảng tứ  phân vị của mẫu số liệu ghép nhóm là $\Delta_Q=Q_3-Q_1=8{,}5$.
	} 
\end{bt}

% nguồn: Lop_12_Dai_so_Chuong_3
\begin{bt}%[2D3H1-3]%[ma:9582f2c6]
	\immini{Đo chiều dài của $80$ con cá, kết quả thu được biểu diễn ở biểu đồ tần số tương đối ghép  nhóm như sau:
	Tìm tứ phân vị thứ nhất  của mẫu số liệu ghép nhóm nói trên.
	}{\resizebox{!}{5cm}{\begin{tikzpicture}[line join=round,line cap=round,>=stealth]
		\draw[->]  (0,0)--(8,0) node[above]{Chiều dài (cm)};
		\draw[->]  (0,0)--(0,6.1) node[right]{Tần số tương đối (\%)};
		 \node[below left] at (0,0) {$O$};
		\foreach  \x/\n in {1/27,2/28,3/29,4/30,5/31,6/32}
		\fill (\x,0) circle (1pt)  node[below]{$\n$};
		\foreach \y  in {1,2,3,4,5}
		 \fill (0,\y) circle (1pt) node [left] {$\y0$};
		 \foreach \x/\y/\p in {1/0.5/5,2/2/20,3/5/50,4/1.5/15,5/1/10}{
			\draw (\x,0)  rectangle (\x+1,\y);
			 \node[above] at (\x+0.5,\y) {$\p\%$};
		} 
	\end{tikzpicture}
	 }
	} 
	\shortans{29} 
	\loigiai{Chuyển biểu  đồ tần số tương đối về dạng ghép nhóm như sau
		\begin{center}
			\begin{tabular}{|l|c|c|c|c|c|}
				\hline 
				Chiều dài & $[27; 28)$ & $[28; 29)$ & $[29; 30)$ &  $[30; 31)$ & $[31; 32)$ \\
				\hline 
				Tần số tương đối & $5\%$ &  $20\%$ & $50\%$ & $15\%$ & $10\%$ \\
				\hline 
			\end{tabular}
		\end{center}
		Ta biết rằng nếu thứ phân vị thứ $k$ là $\dfrac{1}{2}\left( x_m+x_{m+1} \right)$,  trong đó $x_m$ và $x_{m+1}$ lần lượt thuộc hai nhóm liên tiếp $\left[u_{j-1};u_j \right)$ và $\left[u_j;u_{j+1} \right)$ thì $Q_k=u_j$.\\
		Áp dụng cho bài trên, tần số tương đối tích lũy $5\%+20\%=25\%$ nên  $Q_1=29$.
	 }
\end{bt}

% nguồn: Lop_12_Dai_so_Chuong_3
\begin{bt}%[2D3H1-3]%[ma:a307d7d3]
	\immini{Điểm thi môn Toán (thang điểm 100, điểm được làm tròn đến 1) của 60 thí sinh được cho mẫu số liệu ghép nhóm như sau 
	Tính khoảng tứ phân vị của mẫu  số liệu ghép nhóm.}
	{\begin{tabular}{|c|c|}
	\hline 
	Số  giờ làm thêm & Số sinh viên \\
	 \hline $[0;9{,}5)$&$1$\\
	\hline $[9{,}5;19{,}5)$&$2$\\ 
	 \hline$[19{,}5;29{,}5)$& $4$\\
	\hline $[29{,}5;39{,}5)$& $6$\\ 
	 \hline $[39{,}5;49{,}5)$& $15$\\
	\hline $[49{,}5;59{,}5)$&  $12$\\
	\hline  $[59{,}5;69{,}5)$& $10$\\
	 \hline $[69{,}5;79{,}5)$& $6$\\
	\hline $[79{,}5;89{,}5)$&  $3$\\
	 \hline $[89{,}5;99{,}5)$& $1$\\
	\hline 
\end{tabular}}
\shortans{23,7} 
 \loigiai{
	Cỡ mẫu là  $n=60$.\\
	Gọi $x_1, x_2, \ldots, x_{60}$ là điểm thi môn Toán của 60 thí sinh và giả sử dãy này đã được sắp xếp  theo thứ tự tăng dần.\\
	Tứ phân  vị thứ nhất $Q_1$ là $\dfrac{x_{15}+x_{16}}{2}$. \\
	Do $x_{15}$ và $x_{16}$ đều thuộc nhóm $[39{,}5;49{,}5)$ nên nhóm này chứa $Q_1$. \\ 
	Do đó, $p=5$; $a_5=39{,}5;$ m $m_5=15$;  $m_1+m_2+m_3+m_4=13$; $a_6-a_5=10$ và ta có
	 \[Q_1=39{,}5+\dfrac{\dfrac{60}{4}-13}{15}\cdot 10\approx 40{,}83.\]
	Tứ  phân vị thứ ba $Q_3$ là $\dfrac{x_{45}+x_{46}}{2}$. \\
	Do $x_{45}$ và $x_{46}$ đều thuộc nhóm $[59{,}5;69{,}5)$ nên nhóm này chứa  $Q_3$.\\
	Do đó $p=7$; $a_7=59{,}5$; $m_7=10$; $m_1+m_2+m_3+m_4+m_5+m_6=40$; $a_6-a_5=10$ và ta  có
	\[Q_3=59{,}5+\dfrac{\dfrac{3\cdot 60}{4}-40}{10}\cdot 10=64{,}5.\] 
	 Khoảng tứ phân vị của mẫu số liệu ghép nhóm là $\Delta_Q=23{,}67$.
 }
\end{bt}

% nguồn: Lop_12_Dai_so_Chuong_3
\begin{bt}%[2D3H1-3]%[ma:29b84ee6]
	Số điểm một  cầu thủ bóng rổ ghi được trong $20$ trận đấu được cho ở bảng sau
	\begin{center}
		\begin{tabular}{|c|c|c|c|c|}
			 \hline
			Điểm số & $[5{,}5;10{,}5)$& $[10{,}5;15{,}5)$& $[15{,}5;20{,}5)$& $[20{,}5;25{,}5)$\\ 
			\hline 
			Số trận & $3$& $9$& $2$& $6$\\ 
			\hline 
		\end{tabular}
	\end{center}
	Hãy tìm khoảng tứ phân vị của  mẫu số liệu ghép nhóm trên.
	 \shortans{9,7}
	\loigiai{ 
		Gọi $x_1;x_2;x_3;\ldots;x_{20}$ lần lượt  là số trận theo thứ tự không giảm.\\
		Ta có 
		\begin{itemize}
			\item $x_1, \ldots, x_3 \in  [5{,}5;10{,}5)$;
			\item $x_4,  \ldots, x_{12} \in [10{,}5;15{,}5)$;
			\item  $x_{13}, x_{14} \in [15{,}5;20{,}5)$;
			\item $x_{15}, \ldots, x_{20} \in [20{,}5;25{,}5)$. 
		\end{itemize}
		Tứ phân vị thứ hai của  dãy số liệu là $\dfrac{1}{2} \left(x_{10}+x_{11} \right)$ thuộc nhóm $[10{,}5;15{,}5)$.\\
		Do đó tứ phân  vị thứ hai của mẫu số liệu là \[Q_2=10{,}5+\dfrac{\dfrac{20}{2}-3}{9} \left(15{,}5-10{,}5\right)=14{,}4.\]
		Tứ phân vị thứ nhất của dãy số liệu  là $\dfrac{1}{2} \left(x_5+x_6 \right)$ thuộc nhóm $[10{,}5;15{,}5)$.\\
		Do đó tứ phân vị thứ nhất của mẫu số liệu  là \[Q_1=10{,}5+\dfrac{\dfrac{20}{4}-3}{9} \left(15{,}5-10{,}5\right)=11{,}6.\]
		Tứ phân vị thứ ba của dãy số liệu là $\dfrac{1}{2} \left(x_{15}+x_{16}  \right)$ thuộc nhóm $[20{,}5;25{,}5)$.\\
		Do đó tứ phân vị thứ  ba của mẫu số liệu là \[Q_3=20{,}5+\dfrac{\dfrac{3.20}{4}-14}{6} \left(25{,}5-20{,}5\right)=21{,}3.\]
		Khoảng tứ phân  vị của mẫu số liệu ghép nhóm là $\Delta_Q=9{,}7$.
	 }
\end{bt}

% nguồn: Lop_12_Dai_so_Chuong_3
\begin{bt}%[2D3H1-3]%[ma:fc705923]
	Thời giảm luyện  tập trong một ngày (tính theo giờ) của một số vận động viên được ghi lại ở bảng sau
	\begin{center}
		\begin{tabular}{|c|c|c|c|c|c|}
			 \hline
			Thời gian luyện tập (giờ) & $[0;2)$&  $[2;4)$& $[4;6)$& $[6;8)$& $[8;10)$\\
			 \hline
			Số vận động viên & $3$& $8$&  $12$& $12$& $4$\\
			 \hline
		\end{tabular}
	\end{center}
	Hãy tìm khoảng tứ  phân vị của mẫu số liệu ghép nhóm trên.
	 \shortans{3,35}
	\loigiai{ 
		Cỡ mẫu  $n=39$.\\
		Tứ phân vị thứ nhất của dãy số liệu $x_1;x_2;\ldots;x_{39}$ là $x_{10} \in [2;4)$.\\ 
		Do đó tứ phân vị thứ  nhất của mẫu số liệu ghép nhóm là \[Q_1=2+\dfrac{\dfrac{1\cdot 39}{4}-3}{8} \left(4-2\right)=\dfrac{59}{16}.\]
		Tứ phân  vị thứ ba của dãy số liệu $x_1;x_2;\ldots;x_{39}$ là $x_{30} \in [6;8)$.\\
		 Do đó tứ phân vị thứ ba của mẫu số liệu ghép nhóm là
		 \[Q_3=6+\dfrac{\dfrac{3\cdot 39}{4}-\left(3+8+12\right)}{12} \left(8-6\right)=\dfrac{169}{24}.\]
		Khoảng tứ  phân vị của mẫu số liệu ghép nhóm là $\Delta_Q=Q_3-Q_1\dfrac{161}{48}\approx3{,}35$.
	} 
\end{bt}

% nguồn: Lop_12_Dai_so_Chuong_3
\begin{bt}%[2D3H1-3]%[ma:ea43de29]
	Bảng sau thống kê cân nặng của $30$ quả đu đủ được lựa chọn ngẫu nhiên sau khi thu hoạch ở vườn nhà Lan. 
	\begin{center}
		\begin{tabular}{|l|c|c|c|c|c|}
			 \hline
			Cân nặng (g) & $[750; 800)$ & $[800; 850)$ & $[850; 900)$ & $[900; 950)$  & $[950; 1000)$ \\
			 \hline
			Số quả bưởi  & $5$& $10$& $5$& $8$& $2$\\
			\hline 
		\end{tabular}
	\end{center}
	Khoảng tứ phân vị của mẫu số liệu ghép nhóm trên  là bao nhiêu? (Làm tròn đến hàng đơn vị).
	 \shortans{103}
	 \loigiai{
		Cỡ  mẫu $n=30$.\\
		Gọi $x_1$; $x_2$; $\ldots$; $x_{30}$ là số cân nặng của $30$ quả đu đủ thu hoạch ở vường nhà Lan được sắp  xếp theo thứ tự không giảm.\\
		Tứ phân vị thứ nhất  là $Q_1=800+\dfrac{\dfrac{30}{4}-5}{10} \cdot 50=812{,}5$.\\
		Tứ  phân vị thứ nhất là $Q_3=900+\dfrac{\dfrac{3\cdot 30}{4}-(5+10+5)}{6}\cdot 50=915{,}625$.\\
		 Khoảng tứ phân vị của mẫu số liệu ghép nhóm trên là $\Delta_Q=Q_3-Q_1=103$.
	 }
\end{bt}

% nguồn: Lop_12_Dai_so_Chuong_3
\begin{bt}%[2D3H1-3]%[ma:39a66a22]
	Một người thống kê lại  thời gian thực hiện các cuộc gọi điện thoại của người đó trong một tuần ở bảng sau
	\begin{center}
		\begin{tabular}{|c|c|c|c|c|c|c|}
			\hline 
			Thời gian cuộc gọi (giây) & $[0;60)$& $[60;120)$& $[120;180)$& $[180;240)$& $[240;300)$&  $[300;360)$\\
			\hline 
			Số  cuộc gọi & $8$& $10$& $7$& $5$& $2$& $1$\\
			 \hline
		\end{tabular}
	\end{center}
	Hãy  tìm khoảng tứ phân vị của mẫu số liệu ghép nhóm trên.
	 \shortans{116}
	\loigiai{ 
		 Số lần thực hiện cuộc gọi là $n=8+10+7+5+2+1=33$.\\
		Gọi $x_1;x_2;x_3;\ldots;x_{33}$ lần lượt  là thời gian thực hiện cuộc gọi theo thứ tự không giảm.\\
		Ta có 
		\begin{itemize}
			\item $x_1, \ldots, x_8  \in [0;60)$;
			\item $x_9, \ldots, x_{18}  \in [60;120)$;
			\item  $x_{19}, \ldots, x_{25} \in [120;180)$;
			\item $x_{26}, \ldots, x_{30} \in  [180;240)$;
			\item $x_{31}, \ x_{32} \in [240;300)$; 
			\item $x_{33} \in  [300;360)$.
		\end{itemize}
		Tứ phân vị thứ nhất của dãy số  liệu là $\dfrac{1}{2} \left(x_8+x_9 \right)$ thuộc nhóm $[60;120)$.\\
		Do đó, tứ phân vị thứ nhất của  mẫu số liệu là \[Q_1=60+\dfrac{\dfrac{33}{4}-8}{10} (120-60)=61{,}5.\]
		Tứ phân vị thứ ba của dãy số liệu là $\dfrac{1}{2} \left(x_{24}+x_{25} \right)$ thuộc nhóm  $\left[120;180\right)$.\\
		Do đó tứ phân vị thứ ba của mẫu số liệu là \[Q_3=120+\dfrac{\dfrac{3\cdot 33}{4}-18}{7}  \left(180-120\right)=\dfrac{1\,245}{7}.\]
		Khoảng tứ  phân vị của mẫu số liệu ghép nhóm là $\Delta_Q=\dfrac{1\,245}{7}-\dfrac{123}{2}=\dfrac{1\,629}{14}\approx 116$.
	 }
\end{bt}

% nguồn: Lop_12_Dai_so_Chuong_3
\begin{bt}%[2D3H1-3]%[ma:adf40acc]
	Thời gian chờ khám bệnh của các bệnh nhân tại phòng khám $X$ được cho trong bảng  sau:
	\begin{center}
		\begin{tabular}{|l|c|c|c|c|}
			\hline Thời gian (phút) & {$[0; 5)$} &  {$[5; 10)$} & {$[10; 15)$} & {$[15; 20)$} \\
			\hline Số bệnh nhân & $3$ & $12$  & $15$ & $8$ \\
			\hline 
		\end{tabular}
	\end{center}
	\shortans{6,79} 
	Tìm khoảng tứ phân vị của mẫu số liệu ghép nhóm này. (Làm tròn đến hàng  phần trăm).
	\loigiai{ 
		Cỡ  mẫu là $n=3+12+15+8=38$.\\
		Gọi $x_1,\ldots,x_{38}$ là thời gian chờ khám bệnh của 38 bệnh nhân này và giả sử rằng dãy số liệu gốc này đã được  sắp xếp theo thứ tự không giảm.\\
		Tứ phân vị thứ nhất của mẫu số liệu gốc  là $x_{10}$ nên nhóm chứa tứ phân vị thứ nhất là nhóm $[5;10)$ và ta có $Q_1=5+\left[\dfrac{\dfrac{38}{4}-3}{12}\right]\cdot 5\approx 7{,}71$.\\
		Tứ phân vị thứ ba của  mẫu số liệu gốc là $x_{29}$ nên nhóm chứa tứ phân vị thứ ba là nhóm $[10;15)$ và ta có $Q_3=10+\left[\dfrac{\dfrac{3\cdot 38}{4}-15}{15}\right]\cdot 5=14{,}5$.\\
	 Vậy khoảng tứ phân vị của mẫu số liệu ghép nhóm là $\Delta_Q=Q_3-Q_1\approx 14{,}5-7{,}71=6{,}79$.}
\end{bt}

% nguồn: Lop_12_Dai_so_Chuong_3
\begin{bt}%[2D3V1-3]%[ma:7e5bbac9]
	Lương tháng của $50$  nhân viên một công ty được biểu diễn ở biểu đồ sau:
	\begin{center}
		\begin{tikzpicture}[scale=1,>=stealth]
				 \pgfplotsset{compat=1.16}
				\begin{axis}[
						 ybar,
						bar width=1, 
						 ymin=0,
						ymax=.5, 
						 enlarge x limits={abs=2},
						xtick=data, 
						height=7cm,width=9.6cm, 
						point meta=y*100, 
						 nodes near coords={\pgfmathprintnumber\pgfplotspointmeta\%},
						nodes  near coords style={/pgf/number format/precision=3,
						/pgf/number format/fixed}, 
						 xticklabels={,,},
						 yticklabel={\pgfmathparse{\tick*100}\pgfmathprintnumber{\pgfmathresult}},
						 tickwidth=0,
						 ymajorgrids,
						axis  y line*=left,
						 axis x line*=left,]
					 \addplot[fill=black!20] table[x expr=\coordindex,y index=0,row sep=\\] {0.08\\0.24\\0.4\\0.16\\0.12\\};
				\end{axis}
				\draw[->,line width=1 pt]  (0:0)--(0:8.5) node[below right,pos=.9,rectangle,minimum width=3.5cm,minimum height=1cm,text width=3.5cm] {Lương (triệu đồng)};
				\draw[->,line width=1 pt] (0:0)--(90:7) node[above left,pos=.8,rectangle,minimum width=2.5cm,minimum height=1cm,text  width=2.5cm] {Tần số tương đối (\%)};
				\foreach \x/\y  in {1.5/6,2.5/8,3.5/10,4.5/12,5.5/14,6.5/16}
				 \draw[thin] (\x,1pt)--(\x,-1pt) node [below] {$\y$};
			\end{tikzpicture}
	\end{center}
	Tính khoảng tứ phân vị của mẫu số liệu ghép nhóm  trên (đơn vị: triệu đồng). Làm tròn kết quả đến hàng phần trăm.
	\shortans{2,96} 
	 \loigiai{
		Ta có  bảng tần số ghép nhóm
		\begin{center}
			\begin{tabular}{|m{3cm}|m{2cm}|m{2cm}|m{2cm}|m{2cm}|m{2cm}|}
				\hline 
				\centering\arraybackslash Nhóm & 
				 \centering\arraybackslash$\left[6;8\right)$ &
				 \centering\arraybackslash$\left[8;10\right)$ &
				 \centering\arraybackslash$\left[10;12\right)$ &
				\centering\arraybackslash$\left[12;14\right)$  &
				\centering\arraybackslash$\left[14;16\right)$  \\
				\hline 
				\centering\arraybackslash Tần số & 
				\centering\arraybackslash$4$  &
				\centering\arraybackslash$12$  &
				\centering\arraybackslash$20$  &
				\centering\arraybackslash$8$  &
				\centering\arraybackslash$6$ \\ 
				 \hline
			\end{tabular}
		\end{center}
		Tứ phân vị thứ  nhất của mẫu số liệu gốc là $x_{13}\in \left[8;10\right)$.\\ Do đó, tứ phân vị thứ nhất của mẫu số liệu ghép nhóm là
		 $$Q_1=8+\dfrac{\dfrac{50}{4}-4}{12}\left(10-8\right)=\dfrac{113}{12}.$$
		Tứ phân vị thứ ba của mẫu số liệu gốc  là $x_{38}\in \left[12;14\right)$.\\ Do đó, tứ phân vị thứ ba của mẫu số liệu ghép nhóm là
		 $$Q_3=12+\dfrac{\dfrac{3 \cdot 50}{4}-36}{8}\left(14-12\right)=\dfrac{99}{8}.$$
		Khoảng tứ phân vị  $\Delta Q=Q_3-Q_1=\dfrac{113}{12}-\dfrac{99}{8}=\dfrac{71}{24}=2{,}96$.
	 }
\end{bt}


%%%%%%%%%%%%%%%% TỰ LUẬN %%%%%%%%%%%%%%%%

% nguồn: Lop_12_Dai_so_Chuong_3
\begin{bt}%[2D3H1-3]%[ma:9817a6bc]
	Hằng ngày ông Thắng đều đi xe buýt từ nhà đến cơ quan. Dưới đây là bản thống kê thời gian của $100$ lần ông Thắng đi xe buýt từ  nhà đến cơ quan
	\begin{longtable}{|c|c|c|c|c|c|c|}
		 \hline
		Thời  gian (phút) & $[15;18)$ & $[18;21)$ & $[21;24)$ & $[24;27)$ & $[27;30)$ & $[30;33)$\\
		\hline 
		Số  lần & $22$ & $38$ & $27$ & $8$ & $4$ & $1$\\
		 \hline
	\end{longtable}
	\begin{enumerate}
		\item Hãy tìm khoảng tứ phân vị của mẫu số liệu ghép nhóm phía trên (làm tròn kết quả đến  hàng phần trăm).
		\item Biết rằng trong $100$ lần đi trên, chỉ có  đúng một lần ông Thắng đi hết $32$ phút. Thời gian của lần đi đó có phải giá trị ngoại lệ không?
	\end{enumerate}
	 \loigiai{
		\begin{enumerate}
			\item Cỡ mẫu $n=100$.  \\
			Gọi $x_1, x_2, x_3, \ldots, x_{100}$ là mẫu số liệu gốc gồm $100$  lần đi xe buýt của ông Thắng. Khi đó ta có
			$x_1,  \ldots, x_{22} \in \left[15; 18\right)$; $ x_{23}, \ldots, x_{60} \in \left[18; 21\right)$; $ x_{61}, \ldots, x_{87} \in \left[21; 24\right)$; $ x_{88}, \ldots., x_{95} \in \left[21; 27\right)$; $ x_{96}, \ldots, x_{99} \in \left[27; 30\right)$; $ x_{100} \in \left[30; 33\right)$.\\
			Tứ phân vị thứ nhất của mẫu số liệu gốc là $\dfrac{1}{2} \left(x_{75}+x_{76} \right)\in  \left[18; 21\right)$.\\
			Do đó tứ phân  vị thứ nhất của mẫu số liệu ghép nhóm là \[Q_1=18+\dfrac{\dfrac{100}{4}-22}{38} \left(21-18\right)=\dfrac{693}{38}.\]
			Tứ phân vị thứ ba của mẫu số liệu gốc  là
			\[\dfrac{1}{2} \left(x_{75}+x_{76} \right)\in  \left[21; 24\right).\]
			Do đó  tứ phân vị thứ ba của mẫu số liệu ghép nhóm là
			 \[Q_3=18+\dfrac{\dfrac{3.100}{4}-\left(22+38\right)}{27} \left(24-21\right)=\dfrac{68}{3}.\]
			Vậy khoảng tứ phân vị của mẫu số liệu ghép nhóm là $\Delta_Q=\dfrac{68}{3}-\dfrac{693}{38}=\dfrac{505}{114} \approx 4{,}43$. 
			\item Trong lần duy nhất ông Thắng đi hết $32$  phút, thời gian đi của ông thuộc nhóm $\left[30; 33\right)$.\\
			Vì $Q_3+1{,}5\Delta_Q \approx  29{,}31< 30$ nên đây là giá trị ngoại lệ của mẫu số liệu ghép nhóm.
		\end{enumerate}
	 }
\end{bt}
